\chapter{Classical Sensor Fusion Techniques}
\label{sec:classicalSF}
%
%
% ============================================

% ============================================

% Springer Nature: each chapter should have a 10--15 line online abstract.
% This chapter abstract has not yet been written in the uploaded manuscript.
\abstract{This chapter reviews the classical, model-based techniques on which sensor fusion rests, and presents them as variations of a single idea: measurements are pieces of information about a hidden state, to be weighted by their uncertainty and combined additively when they are independent. The chapter starts from the static inverse problem and develops ordinary, weighted, regularised and sequential least squares, their maximum-likelihood and maximum-a-posteriori interpretations, the information matrix, and Gauss--Newton iterations for nonlinear sensor models. The same prediction--correction structure is then extended to dynamic states through Bayesian filtering, covering the Kalman filter, the extended and unscented Kalman filters, particle filters, innovation-based gating and the normalised innovation squared as a consistency diagnostic; a worked automotive-radar cut-in example with real numbers accompanies the discussion. Set-membership and zonotopic estimation are presented as the bounded-error counterpart of probabilistic filtering, with strip intersections, order reduction and a direct comparison with the EKF on the same scenario. Track-level fusion and data association follow, including temporal and spatial alignment, gating, nearest-neighbour, probabilistic and multi-hypothesis association, track management, fusion with known cross-correlation and covariance intersection. Evidence theory (Dempster--Shafer, the transferable belief model and DSmT) is introduced for conflicting qualitative information, optimisation-based fusion and factor graphs for trajectory estimation, and multi-criteria decision-making, in particular incomplete AHP solved by logarithmic least squares, for context-dependent sensor reliability weighting. The chapter closes with ten problems, whose solutions are collected at the end of the book and in the companion repository.}
\suppressfloats[t]

\section{Mathematical introduction to sensor fusion}
\label{sec:math-intro}
% ============================================

Sensor fusion concerns the estimation of an unknown physical quantity from multiple measurements that may differ in modality, dimension, accuracy, sampling rate, and uncertainty. The quantity to be estimated may be a static parameter, such as the location of a stationary source, or a time-varying state, such as the position, velocity, and orientation of a moving vehicle. Table~\ref{tab:notation} collects the notation used throughout the book and records where symbols are reused with different meanings across the method families reviewed here.
\begin{table}[!t]
    \centering
    \small
    \caption{Principal notation. Symbols are reused across method families where the literatures conflict; the scope column states where each reading applies.}
    \label{tab:notation}
    \begin{tabular}{p{0.13\linewidth} p{0.44\linewidth} p{0.28\linewidth}}
        \hline
        \textbf{Symbol} & \textbf{Meaning} & \textbf{Scope} \\
        \hline
        $\mathbf{x}$, $\widehat{\mathbf{x}}$ & state, estimate & throughout \\
        $\mathbf{y}$, $\mathbf{r}$ & measurement, measurement noise & throughout \\
        $\mathbf{g}(\cdot)$, $\mathbf{G}$ & measurement function, measurement matrix & throughout \\
        $\mathbf{f}(\cdot)$, $\mathbf{q}$ & process model, process noise & dynamic models \\
        $\mathbf{R}$, $\mathbf{Q}$, $\mathbf{P}$ & measurement, process, state covariance & probabilistic fusion \\
        $\mathbf{\Lambda}$ & information matrix, $\mathbf{P}^{-1}$ & probabilistic fusion \\
        $\boldsymbol{\nu}$, $\mathbf{S}$, $\mathbf{K}$ & innovation, its covariance, gain & filtering \\
        $\mathbf{m}$ & prior mean & regularised least squares \\
        $p$, $H$, $\gamma$ & zonotope centre, generators, strip gain & set-membership \\
        $\Omega$, $m(\cdot)$, $K$ & frame of discernment, belief mass, conflict & evidence theory \\
        $w$, $q$, $\mu$ & fusion weight, quality indicator, membership & reliability-aware fusion \\
        $\mathcal{G}$, $\mathcal{N}_i$ & communication graph, neighbour set & multi-agent fusion \\
        $\mathbf{H}$ & measurement Jacobian or matrix inside filtering recursions (equals $\mathbf{G}$ for linear models) & Kalman filtering, Chapters~\ref{sec:learning}--\ref{sec:applications} \\
        $\mathbf{L}$, $\lambda_2$, $\epsilon$ & graph Laplacian, algebraic connectivity, consensus step size & multi-agent fusion \\
        $\alpha_i$ & discount (reliability) factor applied to the evidence of source $i$ & evidence theory, reliability-aware fusion \\
        NIS, NEES & normalised innovation squared, normalised estimation error squared (consistency statistics) & filtering, validation \\
        ECE, $T$ & expected calibration error, temperature-scaling parameter & learning-based fusion \\
        $\mathbf{f}_c$, $\mathbf{f}_l$, $\mathbf{f}_r$ & camera, LiDAR, radar feature representations & deep multimodal fusion \\
        $\mathbf{T}_{VS}$, $\mathbf{K}$ & extrinsic sensor-to-vehicle transform in $SE(3)$; camera intrinsic matrix (the gain is written $\mathbf{K}_k$ when both appear) & calibration and alignment, Chapter~\ref{sec:application} \\
        \added{$\mathbf{\Omega}_i$} & \added{information matrix of the $i$-th factor in factor-graph and SLAM formulations (bold, unrelated to the frame of discernment $\Omega$)} & \added{factor graphs, Chapter~\ref{sec:learning}} \\
        \hline
    \end{tabular}
\end{table}

\begin{addedblock}
Readers trained in machine learning or computer science will meet most objects of this chapter under other names, because the estimation and the machine-learning literatures describe the same probabilistic models from different starting points \cite{bishop2006_prml,koller2009_pgm,sarkka2013_bayesian}. Table~\ref{tab:ml-dictionary} gives the correspondence for the main concepts of this and the following chapters. Some entries are exact identities: the Kalman filter is exact inference in a linear-Gaussian state-space model, that is, in a hidden Markov model with continuous states, and the regularised least-squares estimate with a Gaussian prior is ridge regression. Others are analogies that help intuition: zonotopes are also an abstract domain for the static analysis of programs and for the verification of neural networks \cite{ghorbal2009_zonotope,gehr2018_ai2}; the logarithmic least-squares completion of an incomplete AHP matrix is a least-squares ranking on a comparison graph, closely related to HodgeRank and to the Bradley--Terry model \cite{jiang2011_hodgerank,bradley1952_paired}; and the resilient consensus of Chapter~\ref{sec:MAS} belongs to the same family as the Byzantine-robust aggregation rules of distributed learning \cite{yin2018_byzantine}.

\begin{table}[!t]
    \centering
    \footnotesize
    \setlength{\tabcolsep}{4pt}
    \caption{The main objects of this book in the vocabulary of machine learning and computer science.}
    \label{tab:ml-dictionary}
    \begin{tabular}{>{\raggedright\arraybackslash}p{0.33\linewidth} >{\raggedright\arraybackslash}p{0.44\linewidth} >{\raggedright\arraybackslash}p{0.14\linewidth}}
        \hline
        \textbf{Estimation and fusion} & \textbf{Machine learning and computer science} & \textbf{In this book} \\
        \hline
        state $\mathbf{x}$ & latent (hidden) variable & Sect.~\ref{sec:math-intro} \\
        measurement model $\mathbf{g}(\cdot)$ with noise $\mathbf{r}$ & observation model, likelihood $p(\mathbf{y}\mid\mathbf{x})$; decoder of a generative model & Sect.~\ref{sec:math-intro} \\
        process model $\mathbf{f}(\cdot)$ with noise $\mathbf{q}$ & transition model $p(\mathbf{x}_k\mid\mathbf{x}_{k-1})$ & Sect.~\ref{sec:bayes-filtering} \\
        (weighted) least squares; maximum-likelihood estimate under Gaussian noise & (weighted) mean-squared-error regression; minimisation of the negative log-likelihood & Sect.~\ref{sec:math-intro} \\
        regularised least squares; MAP estimate with Gaussian prior & ridge (Tikhonov, $\ell_2$) regularisation, the prior acting as regulariser & Sect.~\ref{sec:math-intro} \\
        information matrix $\mathbf{\Lambda}$ & Fisher information; Hessian of the negative log-likelihood & Sect.~\ref{sec:math-intro} \\
        sequential least squares & recursive least squares; online learning & Sect.~\ref{sec:math-intro} \\
        Bayesian filter (prediction and update) & forward recursion of inference in a hidden Markov model or dynamic Bayesian network & Sect.~\ref{sec:bayes-filtering} \\
        Kalman filter & exact inference in a linear-Gaussian state-space model (linear dynamical system) & Sect.~\ref{sec:bayes-filtering} \\
        EKF, UKF & approximate inference by linearisation and by deterministic sampling & Sect.~\ref{sec:bayes-filtering} \\
        particle filter & sequential Monte Carlo (sequential importance resampling) & Sect.~\ref{sec:bayes-filtering} \\
        innovation $\boldsymbol{\nu}$, NIS test & prediction residual; calibration and out-of-distribution test on residuals & Sect.~\ref{sec:bayes-filtering} \\
        zonotopic set-membership estimate & reachability analysis; zonotope abstract domain of program analysis and neural-network verification & Sect.~\ref{sec:set-membership} \\
        gating and data association & outlier rejection; bipartite matching (assignment problem) & Sect.~\ref{sec:track-level} \\
        covariance intersection & conservative aggregation of estimates with unknown dependence, avoiding double counting & Sect.~\ref{subsec:track-fusion} \\
        belief masses, Dempster's rule & belief functions; second-order uncertainty in classification, as in evidential deep learning \cite{sensoy2018_edl} & Sect.~\ref{sec:evidence} \\
        factor graph & probabilistic graphical model \cite{kschischang2001_factor}; sparse nonlinear least squares & Sect.~\ref{sec:factor-graphs} \\
        incomplete AHP, logarithmic least squares & ranking from pairwise comparisons (HodgeRank, Bradley--Terry) & Sect.~\ref{sec:mcdm} \\
        learned measurement covariance & heteroscedastic regression; uncertainty calibration & Chapter~\ref{sec:learning} \\
        consensus; resilient (W-MSR) consensus & gossip averaging \cite{boyd2006_gossip}; Byzantine-robust aggregation & Chapter~\ref{sec:MAS} \\
        \hline
    \end{tabular}
\end{table}
\end{addedblock}






% ============================================
\subsection{Sensor fusion as an inverse problem}
% ============================================
Let $\mathbf{x}\in\mathbb{R}^{d_x}$ denote the unknown state or parameter vector. The $n$-th sensor measurement is modelled as
\begin{align}
\mathbf{y}_n=\mathbf{g}_n(\mathbf{x})+\mathbf{r}_n, \qquad n=1,\ldots,N,
\end{align}
where $\mathbf{y}_n\in\mathbb{R}^{d_n}$ is the observed measurement, $\mathbf{g}_n:\mathbb{R}^{d_x}\rightarrow\mathbb{R}^{d_n}$ is the measurement or observation function of sensor $n$, and $\mathbf{r}_n\in\mathbb{R}^{d_n}$ represents measurement noise and modelling error.

The sensor-fusion problem is therefore an inverse problem: the forward models $\mathbf{g}_n(\cdot)$ describe how a hypothetical state produces measurements, while fusion attempts to infer the state that most plausibly explains the measurements actually observed. 

Different sensors need not measure the state directly. For example, a GNSS receiver may provide a noisy position measurement,
\begin{align*}
\mathbf{y}_{\mathrm{GNSS}}=\mathbf{p}+\mathbf{r}_{\mathrm{GNSS}},
\end{align*}
whereas a range sensor located at $\mathbf{s}_n$ may provide
\begin{align*}
y_n=\left\|\mathbf{p}-\mathbf{s}_n\right\|_2+r_n.
\end{align*}
A camera may instead provide a perspective projection,
\begin{align*}
\Pi\left(\boldsymbol{\mathcal{R}}\,\mathbf{p}+\mathbf{t}\right) +\mathbf{r}_n,
\end{align*}
where $\Pi(\cdot)$ is a nonlinear projection operator and $\boldsymbol{\mathcal{R}}$ is a rotation matrix. Sensor fusion places all such measurements within a common estimation framework despite their different physical interpretations.

% ============================================
\subsection{Measurement uncertainty}
% ============================================
Measurements are not exact. Their uncertainty can arise from thermal noise, quantisation, imperfect calibration, environmental interference, finite resolution, unmodelled dynamics, or errors in the assumed sensor model.

A standard starting assumption is that the measurement error has zero mean, $\mathbb{E}[\mathbf{r}_n]=\mathbf{0}_{d_n}$, and covariance $\mathbf{R}_n\in\mathbb{R}^{d_n\times d_n}$. The covariance matrix $\mathbf{R}_n$ describes both the magnitude and the directional structure of the uncertainty. For a scalar measurement, $y_n=g_n(\mathbf{x})+r_n$, the covariance reduces to the variance $\operatorname{var}(r_n)=\sigma_n^2$. A larger $\sigma_n^2$ means that the measurement is less precise. For a vector measurement $\mathbf{y}\triangleq \begin{bmatrix} y_{1} & y_{2} & \cdots & y_{d_n} \end{bmatrix}^T \in \mathbb{R}^{d_n}$,
\begin{align*}
\mathbf{R}=
\begin{bmatrix}
\operatorname{var}(r_{1}) & \operatorname{cov}(r_{1},r_{2}) & \cdots & \operatorname{cov}(r_{1},r_{d_n})\\
\operatorname{cov}(r_{2},r_{1}) & \operatorname{var}(r_{2}) & \cdots & \operatorname{cov}(r_{2},r_{d_n}) \\
\vdots & \vdots & \ddots & \vdots\\
\operatorname{cov}(r_{d_n},r_{1})  & \operatorname{cov}(r_{d_n},r_{2}) & \cdots & \operatorname{var}(r_{d_n}) \\
\end{bmatrix},
\end{align*}
the diagonal entries represent component-wise variances, while the off-diagonal entries represent correlations between measurement components. When the off-diagonal terms are non-zero, the measurement components are correlated, and the uncertainty region is tilted with respect to the coordinate axes. An eigenvalue decomposition removes the tilt by rotating into the axes along which the errors are uncorrelated. For a two-dimensional position measurement
the correlated uncertainty region then becomes an ellipse,
%Geometrically, this means the shape of the uncertainty is tilted relative to our standard axes. To simplify this shape, we use eigenvalue decomposition. This tool rotates our coordinate system to align perfectly with the orientation of the uncertainty. In this new coordinate system, the correlation disappears, and the errors become completely independent. For a two-dimensional position measurement, this transformation maps the correlated uncertainty region into an uncertainty ellipse defined by:
\begin{align*}
\mathbf{R}=\mathbf{U}\left[\begin{matrix}\lambda _{1}&0\\ 0&\lambda _{2}\end{matrix}\right]\mathbf{U}^{\mathsf{T}} \;.
\end{align*}
In this decoupled form, the eigenvectors in $\mathbf{U}$ serve as a rotation matrix that determines the principal axes (directions) of the uncertainty ellipse. Meanwhile, the eigenvalues $\lambda _{1}$ and $\lambda _{2}$ represent the decoupled variances, directly dictating the lengths (magnitudes) of the ellipse's semi-axes along those new directions.

% ============================================
\subsection{Independent and correlated measurement errors}
% ============================================
A common simplifying assumption is that the errors of different measurements are mutually independent. Under this assumption, measurements may have different covariance matrices, but their errors do not share common random components. Independence is reasonable when measurements arise from physically separate noise processes. It may be violated, however, when sensors share calibration parameters, environmental disturbances, timing errors, communication links, or preprocessing algorithms. For example, two position estimates derived from the same camera images generally have correlated errors. Likewise, repeated measurements from the same biased sensor may be strongly correlated even when their high-frequency noise samples are independent. The distinction is important because treating correlated measurements as independent effectively counts shared information more than once, causing an estimator to become overconfident\footnote{The intuition is simple: if two sensors share a common disturbance, calibration error, map error, timing error, or modelling error, then their mistakes tend to point in the same direction. Seeing both sensors agree does not mean the estimate is twice as reliable. They may agree because they are affected by the same hidden error.}.

Mathematically, the absence of correlation between two different measurement errors can be described through their cross-covariance. For measurements indexed by $m$ and $n$, this cross-covariance is defined as
\begin{equation}
\operatorname{cov}(\mathbf{r}_m,\mathbf{r}_n) = \mathbb{E}\left[\left(\mathbf{r}_m-\mathbb{E}[\mathbf{r}_m]\right)\left(\mathbf{r}_n-\mathbb{E}[\mathbf{r}_n]\right)^{\mathsf T}\right].
\end{equation}
Thus, the assumption that different measurement errors are uncorrelated is written as
\begin{equation}
\operatorname{cov}(\mathbf{r}_m,\mathbf{r}_n) = \mathbf{0}, \quad m\neq n.
\end{equation}
It is useful to distinguish between being uncorrelated and being independent. Independence implies zero cross-covariance, provided the covariance exists. The reverse is not true in general. However, when the errors are jointly Gaussian, zero cross-covariance is sufficient to imply independence. This is one reason why Gaussian models are especially convenient in sensor fusion.

% ============================================
\subsection{Stacked measurement model}
% ============================================
Sensor fusion becomes particularly transparent when all measurements are written in a common stacked form. Starting from the individual measurement models
\begin{equation}
\mathbf{y}_n
=
\mathbf{g}_n(\mathbf{x})
+
\mathbf{r}_n,
\qquad n=1,\ldots,N,
\end{equation}
we define the stacked measurement vector as
\begin{equation}
\underbrace{\begin{bmatrix}
\mathbf{y}_1\\
\mathbf{y}_2\\
\vdots\\
\mathbf{y}_N
\end{bmatrix}}_{\triangleq \mathbf{y}} = 
\underbrace{\begin{bmatrix}
\mathbf{g}_1(\mathbf{x})\\
\mathbf{g}_2(\mathbf{x})\\
\vdots\\
\mathbf{g}_N(\mathbf{x})
\end{bmatrix}}_{\triangleq \mathbf{g}(\mathbf{x})} + 
\underbrace{\begin{bmatrix}
\mathbf{r}_1\\
\mathbf{r}_2\\
\vdots\\
\mathbf{r}_N
\end{bmatrix}.}_{\triangleq \mathbf{r}}
\end{equation}
This compact representation is important because it allows heterogeneous measurements to be handled in a unified way. The entries of $\mathbf{y}$ may include positions, velocities, accelerations, ranges, bearings, pixel coordinates, signal strengths, or other sensor outputs. The role of $\mathbf{g}(\mathbf{x})$ is to predict, for a given candidate state $\mathbf{x}$, what each sensor should observe.

The covariance of the stacked noise vector is
\begin{equation}
\mathbf{R}
\triangleq
\operatorname{cov}(\mathbf{r}) = \mathbb{E} \left[ \left( \mathbf{r}-\mathbb{E}[\mathbf{r}] \right) \left( \mathbf{r}-\mathbb{E}[\mathbf{r}] \right)^{\mathsf T} \right].
\end{equation}
If the measurement errors are mutually uncorrelated, then the covariance matrix has the block-diagonal structure
\begin{equation}
\mathbf{R}=
\begin{bmatrix}
\mathbf{R}_1 & \mathbf{0} & \cdots & \mathbf{0}\\
\mathbf{0} & \mathbf{R}_2 & \cdots & \mathbf{0}\\
\vdots & \vdots & \ddots & \vdots\\
\mathbf{0} & \mathbf{0} & \cdots & \mathbf{R}_N
\end{bmatrix}.
\end{equation}
If the measurement errors are correlated, the covariance matrix instead contains nonzero off-diagonal blocks:
\begin{equation}
\mathbf{R}
=
\begin{bmatrix}
\mathbf{R}_{11} & \mathbf{R}_{12} & \cdots & \mathbf{R}_{1N}\\
\mathbf{R}_{21} & \mathbf{R}_{22} & \cdots & \mathbf{R}_{2N}\\
\vdots & \vdots & \ddots & \vdots\\
\mathbf{R}_{N1} & \mathbf{R}_{N2} & \cdots & \mathbf{R}_{NN}
\end{bmatrix}.
\end{equation}
The diagonal block $\mathbf{R}_{nn}$ describes the uncertainty of the $n$th measurement, while the off-diagonal block $\mathbf{R}_{mn}$ is
\begin{equation}
\mathbf{R}_{mn} = \operatorname{cov}(\mathbf{r}_m,\mathbf{r}_n).
\end{equation}
These off-diagonal terms should not be ignored when sensors share common sources of uncertainty. If they are ignored, the estimator may behave as if more independent information is available than there really is.

% ============================================
\subsection{Gaussian measurement model}
% ============================================
A common and analytically convenient assumption is that the stacked measurement noise is Gaussian, i.e., $\mathbf{r} \sim \mathcal{N} \left( \mathbf{0}, \mathbf{R} \right)$. Since the measurement model is $\mathbf{y}=\mathbf{g}(\mathbf{x})+\mathbf{r}$, the conditional distribution of the measurements given the state is $\mathbf{y}\mid\mathbf{x} \sim \mathcal{N} \left( \mathbf{g}(\mathbf{x}), \mathbf{R} \right)$. 
The corresponding likelihood is
\begin{align}
\label{eq:Gaussianlikelihood}
p(\mathbf{y}\mid\mathbf{x}) =& \frac{1} {(2\pi)^{d_y/2}|\mathbf{R}|^{1/2}}\times  \\ 
&\exp \left( - \frac{1}{2} \left( \mathbf{y}-\mathbf{g}(\mathbf{x})\right)^{\mathsf T}\mathbf{R}^{-1}\left(\mathbf{y}-\mathbf{g}(\mathbf{x})\right)\right), \nonumber
\end{align}
where $d_y$ is the dimension of the stacked measurement vector. This expression shows that the plausibility of a state is not determined by the residual alone, but by the residual relative to the uncertainty. A residual of a given numerical size may be significant for a precise sensor and insignificant for a noisy sensor.

% ============================================
\subsection{Residuals, least squares, and optimality criteria}
% ============================================
For a candidate state $\mathbf{x}$, the residual is defined as
\begin{align}
\mathbf{e}(\mathbf{x}) \triangleq \mathbf{y} - \mathbf{g}(\mathbf{x}).
\end{align}
The residual measures the discrepancy between the actual measurements and the measurements predicted by the model. Sensor fusion can therefore be formulated as the problem of finding the state that makes the residual small in an appropriate sense.

An estimator is obtained by minimising a cost function:
\begin{align}
\widehat{\mathbf{x}} = \operatorname*{arg\,min}_{\mathbf{x}} J(\mathbf{x}).
\end{align}
Different choices of $J(\mathbf{x})$ correspond to different assumptions about the measurement errors and different notions of optimality. Least-squares methods use a quadratic cost. For scalar measurements, the ordinary least-squares cost is
\begin{align}
J_{\mathrm{LS}}(\mathbf{x}) = \sum_{n=1}^{N}\left(y_n-g_n(\mathbf{x})\right)^2.
\end{align}
For vector-valued measurements, this becomes
\begin{align}
J_{\mathrm{LS}}(\mathbf{x}) &=\sum_{n=1}^{N}\left(\mathbf{y}_n-\mathbf{g}_n(\mathbf{x})\right)^{\mathsf T}\left(\mathbf{y}_n-\mathbf{g}_n(\mathbf{x})\right) \nonumber \\
&= \left( \mathbf{y} - \mathbf{g}(\mathbf{x}) \right)^{\mathsf T}\left(\mathbf{y}-\mathbf{g}(\mathbf{x})\right) \nonumber \\
& = \left\|\mathbf{y} - \mathbf{g}(\mathbf{x}) \right\|_2^2 . % [July 2026 pass] removed stray . left over from a first-round edit
\end{align}
The corresponding ordinary least-squares estimate is
\begin{align}
\widehat{\mathbf{x}}_{\mathrm{LS}} = \operatorname*{arg\,min}_{\mathbf{x}} \left\| \mathbf{y} - \mathbf{g}(\mathbf{x}) \right\|_2^2.
\end{align}
Ordinary least squares is simple and often effective, but it treats all residual components equally. This can be inappropriate in sensor fusion because measurements may have different units, different noise levels, and different reliability.

Consider first the linear case, in which every measurement function is linear, so that the stacked model of the previous subsection reduces to
\begin{align}
\mathbf{y}=\mathbf{G}\mathbf{x}+\mathbf{r},
\label{eq:linear-measurement-model}
\end{align}
where $\mathbf{G}\in\mathbb{R}^{d_y\times d_x}$ in~\eqref{eq:linear-measurement-model} stacks the individual measurement matrices $\mathbf{G}_n$ vertically. Each row of $\mathbf{G}$ states how one scalar measurement responds to the unknown state.
For this linear model, the ordinary least-squares cost is
\begin{align}
J_{\mathrm{LS}}(\mathbf{x}) = \left( \mathbf{y} - \mathbf{G}\mathbf{x} \right)^{\mathsf T} \left( \mathbf{y} - \mathbf{G}\mathbf{x} \right).
\end{align}
The gradient with respect to $\mathbf{x}$ is
\begin{align}
\nabla_{\mathbf{x}} J_{\mathrm{LS}}(\mathbf{x}) = -2 \mathbf{G}^{\mathsf T} \mathbf{y} + 2 \mathbf{G}^{\mathsf T} \mathbf{G} \mathbf{x}.
\end{align}
At the minimum, the gradient is zero. Hence,
\begin{align}
\mathbf{G}^{\mathsf T} \mathbf{G} \widehat{\mathbf{x}}_{\mathrm{LS}} = \mathbf{G}^{\mathsf T} \mathbf{y}.
\end{align}
If $\mathbf{G}^{\mathsf T}\mathbf{G}$ is invertible, the solution is
\begin{align}
\widehat{\mathbf{x}}_{\mathrm{LS}} = \left( \mathbf{G}^{\mathsf T} \mathbf{G} \right)^{-1} \mathbf{G}^{\mathsf T} \mathbf{y}.
\end{align}
The invertibility of $\mathbf{G}^{\mathsf T}\mathbf{G}$ requires the columns of $\mathbf{G}$ to be linearly independent. In sensor-fusion terms, this means that the available measurements must contain enough independent information to identify all components of $\mathbf{x}$. If this condition fails, the estimate may be non-unique or highly sensitive to noise.
Substituting the model $\mathbf{y}=\mathbf{G}\mathbf{x}+\mathbf{r}$ into the estimator gives
\begin{align}
\widehat{\mathbf{x}}_{\mathrm{LS}} = \mathbf{x} + \left( \mathbf{G}^{\mathsf T} \mathbf{G} \right)^{-1} \mathbf{G}^{\mathsf T} \mathbf{r}.
\end{align}
Assuming the noise is zero mean, taking expectations yields
\begin{align}
\mathbb{E} \left[ \widehat{\mathbf{x}}_{\mathrm{LS}} \right] = \mathbf{x},
\end{align}
so the estimator is unbiased when the model is correct and the noise is zero mean. Its covariance is
\begin{align}
\operatorname{cov} \left( \widehat{\mathbf{x}}_{\mathrm{LS}} \right) =
\left(
\mathbf{G}^{\mathsf T}
\mathbf{G}
\right)^{-1}
\mathbf{G}^{\mathsf T}
\mathbf{R}
\mathbf{G}
\left(
\mathbf{G}^{\mathsf T}
\mathbf{G}
\right)^{-1}.
\end{align}
If the noise covariance is $\mathbf{R}=\sigma^2\mathbf{I}$, this simplifies to
\begin{align}
\operatorname{cov} \left(\widehat{\mathbf{x}}_{\mathrm{LS}} \right) = \sigma^2\left( \mathbf{G}^{\mathsf T} \mathbf{G} \right)^{-1}.
\end{align}
This expression shows that estimation uncertainty depends not only on the sensor noise level, but also on the geometry and conditioning of the measurement matrix.

\subsection{Weighted least squares and the information matrix}

Weighted least squares (WLS) incorporates measurement uncertainty directly into the cost function. Instead of penalising all residuals equally, it weights the residual according to the inverse covariance matrix:
\begin{align}
J_{\mathrm{WLS}}(\mathbf{x}) = \left( \mathbf{y} - \mathbf{G}\mathbf{x} \right)^{\mathsf T} \mathbf{R}^{-1} \left( \mathbf{y} - \mathbf{G}\mathbf{x} \right).
\end{align}
This quantity is the squared Mahalanobis norm of the residual:
\begin{align}
J_{\mathrm{WLS}}(\mathbf{x}) = \left\| \mathbf{y} - \mathbf{G}\mathbf{x} \right\|_{\mathbf{R}^{-1}}^2.
\end{align}
The interpretation is as follows: residuals are measured relative to their uncertainty. A residual from a precise sensor is penalised more strongly than the same residual from a noisy sensor. 
Proceeding exactly as in the unweighted case, the gradient
\begin{align}
\nabla_{\mathbf{x}}
J_{\mathrm{WLS}}(\mathbf{x}) = -2 \mathbf{G}^{\mathsf T} \mathbf{R}^{-1} \mathbf{y} + 2 \mathbf{G}^{\mathsf T} \mathbf{R}^{-1} \mathbf{G} \mathbf{x}
\end{align}
vanishes where $$\mathbf{G}^{\mathsf T}\mathbf{R}^{-1}\mathbf{G}\,\widehat{\mathbf{x}}_{\mathrm{WLS}}=\mathbf{G}^{\mathsf T}\mathbf{R}^{-1}\mathbf{y},$$ so that whenever $\mathbf{G}^{\mathsf T}\mathbf{R}^{-1}\mathbf{G}$ is invertible
%The gradient of the WLS cost is
%\begin{align}
%\nabla_{\mathbf{x}}
%J_{\mathrm{WLS}}(\mathbf{x}) = -2 \mathbf{G}^{\mathsf T} \mathbf{R}^{-1} \mathbf{y} %+ 2 \mathbf{G}^{\mathsf T} \mathbf{R}^{-1} \mathbf{G} \mathbf{x}.
%\end{align}
%Setting the gradient to zero gives:
%\begin{align}
%\mathbf{G}^{\mathsf T}
%\mathbf{R}^{-1}
%\mathbf{G}
%\widehat{\mathbf{x}}_{\mathrm{WLS}}
%=
%\mathbf{G}^{\mathsf T}
%\mathbf{R}^{-1}
%\mathbf{y}.
%\end{align}
%If $\mathbf{G}^{\mathsf T}\mathbf{R}^{-1}\mathbf{G}$ is invertible, then
\begin{align}
\widehat{\mathbf{x}}_{\mathrm{WLS}}
=
\left(
\mathbf{G}^{\mathsf T}
\mathbf{R}^{-1}
\mathbf{G}
\right)^{-1}
\mathbf{G}^{\mathsf T}
\mathbf{R}^{-1}
\mathbf{y}.
\end{align}
Substituting $\mathbf{y}=\mathbf{G}\mathbf{x}+\mathbf{r}$ gives
\begin{align}
\widehat{\mathbf{x}}_{\mathrm{WLS}}
=
\mathbf{x}
+
\left(
\mathbf{G}^{\mathsf T}
\mathbf{R}^{-1}
\mathbf{G}
\right)^{-1}
\mathbf{G}^{\mathsf T}
\mathbf{R}^{-1}
\mathbf{r}.
\end{align}
Therefore, for $\mathbb{E}
\left[\mathbf{r}\right]
=\mathbf{0}$, the estimator is unbiased, i.e., $\mathbb{E} \left[\widehat{\mathbf{x}}_{\mathrm{WLS}} \right] = \mathbf{x}$. The covariance of the WLS estimator is
\begin{align}
\operatorname{cov}
\left(
\widehat{\mathbf{x}}_{\mathrm{WLS}}
\right)
=
\left(
\mathbf{G}^{\mathsf T}
\mathbf{R}^{-1}
\mathbf{G}
\right)^{-1}.
\end{align}
The matrix
\begin{align}
\mathbf{\Lambda}
=
\mathbf{G}^{\mathsf T}
\mathbf{R}^{-1}
\mathbf{G}
\end{align}
is often called the \emph{information matrix}. It measures how much information the measurements provide about the unknown state. The covariance is the inverse of this information matrix, i.e., $\operatorname{cov} \left( \widehat{\mathbf{x}}_{\mathrm{WLS}} \right) = \mathbf{\Lambda}^{-1}$.

Weighted least squares has a direct maximum-likelihood interpretation. Assume that $\mathbf{r}
\sim
\mathcal{N}
\left(
\mathbf{0},
\mathbf{R}
\right)$. Then, $\mathbf{y}\mid\mathbf{x}
\sim
\mathcal{N}
\left(
\mathbf{g}(\mathbf{x}),
\mathbf{R}
\right)$.
The maximum-likelihood estimate is defined as
\begin{equation}
\widehat{\mathbf{x}}_{\mathrm{ML}}
=
\operatorname*{arg\,max}_{\mathbf{x}}
p(\mathbf{y}\mid\mathbf{x}).
\end{equation}
Equivalently, it minimises the negative log-likelihood. Ignoring terms that do not depend on $\mathbf{x}$, from~\eqref{eq:Gaussianlikelihood} we obtain
\begin{align*}
-\log p(\mathbf{y}\mid\mathbf{x}) =
\frac{1}{2}
\left(\mathbf{y} - \mathbf{g}(\mathbf{x}) \right)^{\mathsf T}
\mathbf{R}^{-1}
\left(
\mathbf{y}
-
\mathbf{g}(\mathbf{x})
\right)
+
C.
\end{align*}
Therefore,
\begin{equation}
\widehat{\mathbf{x}}_{\mathrm{ML}}
=
\operatorname*{arg\,min}_{\mathbf{x}}
\left(
\mathbf{y}
-
\mathbf{g}(\mathbf{x})
\right)^{\mathsf T}
\mathbf{R}^{-1}
\left(
\mathbf{y}
-
\mathbf{g}(\mathbf{x})
\right).
\end{equation}
Thus, under additive Gaussian measurement noise, $\widehat{\mathbf{x}}_{\mathrm{ML}} = \widehat{\mathbf{x}}_{\mathrm{WLS}}$. This result explains why least-squares methods are so widely used in sensor fusion. They can be derived either from an optimisation viewpoint, by minimising weighted residuals, or from a probabilistic viewpoint, by maximising a Gaussian likelihood.

\subsection{Prior information: regularised least squares and the MAP estimate}
In some problems, measurements alone are insufficient to estimate $\mathbf{x}$ reliably. This may occur when the number of measurements is small, when the measurement geometry is poor, when the inverse problem is ill-conditioned, or when some components of the state are weakly observable. Regularised least squares addresses this by incorporating prior information. Let $\mathbf{m}$ denote a prior estimate of $\mathbf{x}$, and let $\mathbf{P}$ denote the covariance associated with this prior. A small $\mathbf{P}$ means that the prior is trusted strongly; a large $\mathbf{P}$ means that the prior is weak. The regularised least-squares cost is
\begin{align}
J_{\mathrm{ReLS}}(\mathbf{x})
=&
\left(
\mathbf{y}
-
\mathbf{G}\mathbf{x}
\right)^{\mathsf T}
\mathbf{R}^{-1}
\left(
\mathbf{y}
-
\mathbf{G}\mathbf{x}
\right) \nonumber \\
&+
\left(
\mathbf{x}
-
\mathbf{m}
\right)^{\mathsf T}
\mathbf{P}^{-1}
\left(
\mathbf{x}
-
\mathbf{m}
\right).
\end{align}
The first term measures consistency with the sensor data. The second term measures consistency with the prior. Differentiating and setting the gradient to zero gives
\begin{equation}
\left(
\mathbf{G}^{\mathsf T}
\mathbf{R}^{-1}
\mathbf{G}
+
\mathbf{P}^{-1}
\right)
\widehat{\mathbf{x}}_{\mathrm{ReLS}}
=
\mathbf{G}^{\mathsf T}
\mathbf{R}^{-1}
\mathbf{y}
+
\mathbf{P}^{-1}
\mathbf{m}.
\end{equation}
Therefore,
\begin{equation}
\widehat{\mathbf{x}}_{\mathrm{ReLS}}
=
\left(
\mathbf{G}^{\mathsf T}
\mathbf{R}^{-1}
\mathbf{G}
+
\mathbf{P}^{-1}
\right)^{-1}
\left(
\mathbf{G}^{\mathsf T}
\mathbf{R}^{-1}
\mathbf{y}
+
\mathbf{P}^{-1}
\mathbf{m}
\right).
\end{equation}
This estimator combines measurement information and prior information. The measurement information matrix is
\begin{equation}
\mathbf{\Lambda}_{\mathrm{meas}}
=
\mathbf{G}^{\mathsf T}
\mathbf{R}^{-1}
\mathbf{G},
\end{equation}
and the prior information matrix is
\begin{equation}
\mathbf{\Lambda}_{\mathrm{prior}}
=
\mathbf{P}^{-1}.
\end{equation}
The posterior information matrix is their sum:
\begin{equation}
\mathbf{\Lambda}_{\mathrm{post}}
=
\mathbf{G}^{\mathsf T}
\mathbf{R}^{-1}
\mathbf{G}
+
\mathbf{P}^{-1}.
\end{equation}
Hence, the posterior covariance is
\begin{equation}
\mathbf{P}^{+}
=
\left(
\mathbf{G}^{\mathsf T}
\mathbf{R}^{-1}
\mathbf{G}
+
\mathbf{P}^{-1}
\right)^{-1}.
\end{equation}
Unlike ordinary LS and WLS, the regularised estimator is generally pulled toward the prior mean $\mathbf{m}$. From a frequentist perspective, this may introduce bias. From a Bayesian perspective, it reflects the deliberate use of prior knowledge.

Regularised least squares has a natural Bayesian interpretation. Suppose that the prior distribution of the unknown state is Gaussian, i.e., $\mathbf{x}
\sim
\mathcal{N}
\left(
\mathbf{m},
\mathbf{P}
\right)$. Suppose also that the measurement model is linear and Gaussian:
\begin{equation}
\mathbf{y}\mid\mathbf{x}
\sim
\mathcal{N}
\left(
\mathbf{G}\mathbf{x},
\mathbf{R}
\right).
\end{equation}
Bayes' rule gives
\begin{equation}
p(\mathbf{x}\mid\mathbf{y})
=
\frac{
p(\mathbf{y}\mid\mathbf{x})p(\mathbf{x})
}{
p(\mathbf{y})
}.
\end{equation}
The maximum a posteriori estimate is
\begin{equation}
\widehat{\mathbf{x}}_{\mathrm{MAP}}
=
\operatorname*{arg\,max}_{\mathbf{x}}
p(\mathbf{x}\mid\mathbf{y}).
\end{equation}
Equivalently, it minimises the negative log-posterior. Ignoring constants independent of $\mathbf{x}$, this gives
\begin{align}
\widehat{\mathbf{x}}_{\mathrm{MAP}}
=
\operatorname*{arg\,min}_{\mathbf{x}}
\bigg[
&\left(
\mathbf{y}
-
\mathbf{G}\mathbf{x}
\right)^{\mathsf T}
\mathbf{R}^{-1}
\left(
\mathbf{y}
-
\mathbf{G}\mathbf{x}
\right) \nonumber \\
&+
\left(
\mathbf{x}
-
\mathbf{m}
\right)^{\mathsf T}
\mathbf{P}^{-1}
\left(
\mathbf{x}
-
\mathbf{m}
\right)
\bigg].
\end{align}
Therefore,
\begin{equation}
\widehat{\mathbf{x}}_{\mathrm{MAP}}
=
\widehat{\mathbf{x}}_{\mathrm{ReLS}}.
\end{equation}
This result is important because it unifies deterministic optimisation and probabilistic inference. The same estimator can be interpreted either as the minimiser of a regularised cost or as the most probable state after combining a Gaussian prior with Gaussian measurements.

% ============================================
\subsection{Sequential least squares}
% ============================================
\label{subsubsec:sequential_least_squares}

In the batch formulation, all measurements are collected first and the estimate is computed by solving one weighted least-squares problem. This is suitable when the complete data set is available in advance. In many sensor-fusion applications, however, measurements arrive sequentially. For example, a mobile robot may receive one range measurement, then a camera feature, then an inertial reading, and so on. Recomputing the full batch solution after every new measurement is possible, but it is often computationally inefficient and unnecessary.

Sequential least squares provides an equivalent recursive way of updating the estimate. Instead of storing and reprocessing all previous measurements, it summarises the information contained in the past measurements by the current estimate and its covariance. When a new measurement arrives, the estimate is corrected using only the new measurement, the previous estimate, and the previous covariance.

Suppose that after processing measurements up to step $n-1$, the available estimate is $\widehat{\mathbf{x}}_{n-1}$, and its uncertainty is described by the covariance matrix $\mathbf{P}_{n-1}$. The covariance $\mathbf{P}_{n-1}$ quantifies how uncertain the estimator is about the current value of the unknown static parameter $\mathbf{x}$. A new measurement arrives according to the linear model
\begin{align*}
\mathbf{y}_n
=
\mathbf{G}_n\mathbf{x}
+
\mathbf{r}_n,
\quad
\mathbb{E}[\mathbf{r}_n]=\mathbf{0},
\quad
\operatorname{cov}(\mathbf{r}_n)=\mathbf{R}_n.
\end{align*}
Here, $\mathbf{G}_n$ is the measurement matrix of the new sensor observation, and $\mathbf{R}_n$ is the covariance of the new measurement error. We assume that the new measurement error is uncorrelated with the errors already incorporated into $\widehat{\mathbf{x}}_{n-1}$. This assumption is what allows information to be added recursively. The previous estimate can be interpreted as a prior estimate for the unknown vector $\mathbf{x}$. Therefore, after receiving $\mathbf{y}_n$, the updated estimate may be obtained by minimising a regularised weighted least-squares cost:
\begin{align}
J_n(\mathbf{x})
=&
\left(
\mathbf{x}
-
\widehat{\mathbf{x}}_{n-1}
\right)^{\mathsf T}
\mathbf{P}_{n-1}^{-1}
\left(
\mathbf{x}
-
\widehat{\mathbf{x}}_{n-1}
\right) \\
&+
\left(
\mathbf{y}_n
-
\mathbf{G}_n\mathbf{x}
\right)^{\mathsf T}
\mathbf{R}_n^{-1}
\left(
\mathbf{y}_n
-
\mathbf{G}_n\mathbf{x}
\right). \nonumber
\end{align}
The first term penalises deviation from the previous estimate, weighted by the inverse of its covariance. If $\mathbf{P}_{n-1}$ is small, the previous estimate is trusted strongly. If $\mathbf{P}_{n-1}$ is large, the previous estimate is uncertain and the new measurement can influence the update more strongly. The second term penalises disagreement with the new measurement, weighted by the inverse measurement covariance $\mathbf{R}_n^{-1}$. To derive the recursive estimator, we differentiate $J_n(\mathbf{x})$ with respect to $\mathbf{x}$ and set the gradient equal to zero. This gives
\begin{equation}
\mathbf{P}_{n-1}^{-1}
\left(
\widehat{\mathbf{x}}_n
-
\widehat{\mathbf{x}}_{n-1}
\right)
-
\mathbf{G}_n^{\mathsf T}
\mathbf{R}_n^{-1}
\left(
\mathbf{y}_n
-
\mathbf{G}_n
\widehat{\mathbf{x}}_n
\right)
=
\mathbf{0}.
\end{equation}
Rearranging terms gives the normal equation
\begin{align*}
&\left(\mathbf{P}_{n-1}^{-1}+\mathbf{G}_n^{\mathsf T}\mathbf{R}_n^{-1}\mathbf{G}_n\right)\widehat{\mathbf{x}}_n\\
&\qquad=
\mathbf{P}_{n-1}^{-1}\widehat{\mathbf{x}}_{n-1}
+\mathbf{G}_n^{\mathsf T}\mathbf{R}_n^{-1}\mathbf{y}_n.
\end{align*}
Hence, the updated estimate can be written in information form as
\begin{equation}
\begin{aligned}
\widehat{\mathbf{x}}_n
&=
\left(
\mathbf{P}_{n-1}^{-1}
+
\mathbf{G}_n^{\mathsf T}
\mathbf{R}_n^{-1}
\mathbf{G}_n
\right)^{-1}\\
&\quad{}\times
\left(
\mathbf{P}_{n-1}^{-1}
\widehat{\mathbf{x}}_{n-1}
+
\mathbf{G}_n^{\mathsf T}
\mathbf{R}_n^{-1}
\mathbf{y}_n
\right).
\end{aligned}
\end{equation}
The corresponding covariance update in information form is
\begin{equation}
\mathbf{P}_n^{-1}
=
\mathbf{P}_{n-1}^{-1}
+
\mathbf{G}_n^{\mathsf T}
\mathbf{R}_n^{-1}
\mathbf{G}_n.
\end{equation}
This equation is central. It says that the posterior information equals the previous information plus the information contributed by the new measurement. The previous information is $\mathbf{P}_{n-1}^{-1}$, while the new measurement information is $\mathbf{G}_n^{\mathsf T}\mathbf{R}_n^{-1}\mathbf{G}_n$. Thus, a precise and informative measurement increases the information matrix significantly, whereas a noisy or weakly informative measurement contributes little.

Although the information-form update is compact, it is often more intuitive to write the same recursion in correction form. Using the matrix inversion lemma, the updated covariance can be written as
\begin{equation}
\mathbf{P}_n
=
\mathbf{P}_{n-1}
-
\mathbf{P}_{n-1}
\mathbf{G}_n^{\mathsf T}
\left(
\mathbf{G}_n
\mathbf{P}_{n-1}
\mathbf{G}_n^{\mathsf T}
+
\mathbf{R}_n
\right)^{-1}
\mathbf{G}_n
\mathbf{P}_{n-1}.
\end{equation}
Define the gain matrix
\begin{equation}
\mathbf{K}_n
=
\mathbf{P}_{n-1}
\mathbf{G}_n^{\mathsf T}
\left(
\mathbf{G}_n
\mathbf{P}_{n-1}
\mathbf{G}_n^{\mathsf T}
+
\mathbf{R}_n
\right)^{-1}.
\end{equation}
Then the covariance update becomes
\begin{equation}
\mathbf{P}_n
=
\mathbf{P}_{n-1}
-
\mathbf{K}_n
\left(
\mathbf{G}_n
\mathbf{P}_{n-1}
\mathbf{G}_n^{\mathsf T}
+
\mathbf{R}_n
\right)
\mathbf{K}_n^{\mathsf T}.
\end{equation}
Equivalently, using the definition of $\mathbf{K}_n$, this can also be written as
\begin{equation}
\mathbf{P}_n
=
\mathbf{P}_{n-1}
-
\mathbf{K}_n
\mathbf{G}_n
\mathbf{P}_{n-1}.
\end{equation}
In numerical implementations, a symmetric Joseph-form covariance update is often preferred:
\begin{equation}
\mathbf{P}_n
=
\left(
\mathbf{I}
-
\mathbf{K}_n
\mathbf{G}_n
\right)
\mathbf{P}_{n-1}
\left(
\mathbf{I}
-
\mathbf{K}_n
\mathbf{G}_n
\right)^{\mathsf T}
+
\mathbf{K}_n
\mathbf{R}_n
\mathbf{K}_n^{\mathsf T}.
\end{equation}
This form is algebraically equivalent under exact arithmetic, but it is often more robust numerically because it better preserves symmetry and positive semidefiniteness of the covariance matrix.

The estimate itself can now be written in correction form:
\begin{equation}
\widehat{\mathbf{x}}_n
=
\widehat{\mathbf{x}}_{n-1}
+
\mathbf{K}_n
\left(
\mathbf{y}_n
-
\mathbf{G}_n
\widehat{\mathbf{x}}_{n-1}
\right).
\end{equation}
The quantity
\begin{equation}
\boldsymbol{\nu}_n
=
\mathbf{y}_n
-
\mathbf{G}_n
\widehat{\mathbf{x}}_{n-1}
\end{equation}
is called the innovation. It is the discrepancy between the actual measurement $\mathbf{y}_n$ and the measurement predicted from the previous estimate, $\mathbf{G}_n\widehat{\mathbf{x}}_{n-1}$. If the innovation is small, the new measurement is consistent with the current estimate. If the innovation is large, the new measurement suggests that the estimate should be corrected.

The covariance of the innovation is
\begin{equation}
\mathbf{S}_n
=
\mathbf{G}_n
\mathbf{P}_{n-1}
\mathbf{G}_n^{\mathsf T}
+
\mathbf{R}_n.
\end{equation}
The first term, $\mathbf{G}_n\mathbf{P}_{n-1}\mathbf{G}_n^{\mathsf T}$, represents the uncertainty in the predicted measurement caused by uncertainty in the previous estimate. The second term, $\mathbf{R}_n$, represents the uncertainty of the new sensor measurement itself. Therefore, $\mathbf{S}_n$ measures the total uncertainty associated with the innovation.

Using $\mathbf{S}_n$, the gain can be written compactly as
\begin{equation}
\mathbf{K}_n
=
\mathbf{P}_{n-1}
\mathbf{G}_n^{\mathsf T}
\mathbf{S}_n^{-1}.
\end{equation}
This expression shows how the gain balances prior uncertainty and measurement uncertainty. If the previous estimate is very uncertain, $\mathbf{P}_{n-1}$ is large, and the gain tends to be larger. The estimator then trusts the new measurement more. If the new measurement is very noisy, $\mathbf{R}_n$ is large, which increases $\mathbf{S}_n$ and reduces the gain. The estimator then changes only slightly.



Sequential least squares therefore has a clear interpretation. The previous estimate summarises all information from past measurements. The new measurement contributes additional information. The innovation tells us how much the new measurement disagrees with the current prediction. The gain determines how strongly the estimate should react to this disagreement. The covariance update quantifies how much uncertainty remains after the new information is incorporated.

For a static unknown $\mathbf{x}$, sequential least squares is algebraically equivalent to batch weighted least squares, provided the same measurements and covariance assumptions are used. Its advantage is computational and conceptual: it processes data recursively and maintains an explicit uncertainty estimate. When a process model is added so that the state evolves over time, the same prediction-correction structure becomes the basis of the Kalman filter~\cite{kalman_1960} and its extensions, which will be discussed later.

% ============================================
\subsection{Nonlinear sensor-fusion models}
% ============================================

Many practical sensors are nonlinear. Range sensors measure Euclidean distance, bearing sensors measure angles, cameras perform perspective projection, and radio-based localisation systems may depend nonlinearly on position and propagation parameters. The general nonlinear model remains
\begin{equation}
\mathbf{y} = \mathbf{g}(\mathbf{x}) + \mathbf{r}.
\end{equation}
The WLS estimate is defined as
\begin{equation}
\widehat{\mathbf{x}}_{\mathrm{WLS}} = \operatorname*{arg\,min}_{\mathbf{x}} \left( \mathbf{y} - \mathbf{g}(\mathbf{x}) \right)^{\mathsf T} \mathbf{R}^{-1} \left( \mathbf{y} - \mathbf{g}(\mathbf{x}) \right).
\end{equation}
Unlike the linear case, this optimisation problem usually does not have a closed-form solution. Iterative numerical methods are therefore used.

Let the Jacobian matrix of the measurement function be
\begin{equation}
\mathbf{G}_{\mathbf{x}}(\mathbf{x}) = \frac{\partial \mathbf{g}(\mathbf{x})}{\partial \mathbf{x}}.
\end{equation}
For the WLS cost, the gradient is
\begin{equation}
\nabla_{\mathbf{x}} J_{\mathrm{WLS}}(\mathbf{x}) = -2 \mathbf{G}_{\mathbf{x}}^{\mathsf T}(\mathbf{x}) \mathbf{R}^{-1} \left(\mathbf{y} - \mathbf{g}(\mathbf{x})\right).
\end{equation}
A gradient-descent update is
\begin{equation}
\widehat{\mathbf{x}}^{(i+1)} = \widehat{\mathbf{x}}^{(i)}+\gamma^{(i)} \mathbf{G}_{\mathbf{x}}^{\mathsf T} \left(\widehat{\mathbf{x}}^{(i)}\right)\mathbf{R}^{-1}\left(\mathbf{y}-\mathbf{g} \left( \widehat{\mathbf{x}}^{(i)} \right)\right),
\end{equation}
where $\gamma^{(i)}$ is a step size. The step size must be chosen carefully: a very small step leads to slow convergence, while a very large step can increase the cost or cause divergence.

Gauss--Newton method often converges faster because it solves a locally linearised least-squares problem at each iteration. Around the current estimate $\widehat{\mathbf{x}}^{(i)}$, the measurement function is approximated by
\begin{equation}
\mathbf{g}(\mathbf{x}) \approx \mathbf{g} \left( \widehat{\mathbf{x}}^{(i)} \right) + \mathbf{G}_{\mathbf{x}} \left( \widehat{\mathbf{x}}^{(i)} \right) \left( \mathbf{x} - \widehat{\mathbf{x}}^{(i)} \right).
\end{equation}
Define the current residual as
\begin{equation}
\mathbf{e}^{(i)} = \mathbf{y} - \mathbf{g} \left( \widehat{\mathbf{x}}^{(i)} \right).
\end{equation}
The Gauss--Newton increment is
\begin{align}
\Delta \mathbf{x}^{(i)} = &\left(
\mathbf{G}_{\mathbf{x}}^{\mathsf T}
\left(
\widehat{\mathbf{x}}^{(i)}
\right)
\mathbf{R}^{-1}
\mathbf{G}_{\mathbf{x}}
\left(
\widehat{\mathbf{x}}^{(i)}
\right)
\right)^{-1} \times \nonumber \\
&\mathbf{G}_{\mathbf{x}}^{\mathsf T}
\left(
\widehat{\mathbf{x}}^{(i)}
\right)
\mathbf{R}^{-1}
\mathbf{e}^{(i)}.
\end{align}
The estimate is then updated as
\begin{equation}
\widehat{\mathbf{x}}^{(i+1)} = \widehat{\mathbf{x}}^{(i)} + \Delta \mathbf{x}^{(i)}.
\end{equation}
Thus, nonlinear sensor fusion can be viewed as a sequence of locally linear WLS problems. At each iteration, the nonlinear model is linearised, a WLS problem is solved, and the estimate is updated. Near the final estimate, the covariance can be approximated using the local linearisation:
\begin{equation}
\operatorname{cov} \left( \widehat{\mathbf{x}}_{\mathrm{WLS}} \right) \approx \left(\mathbf{G}_{\mathbf{x}}^{\mathsf T}\left(\widehat{\mathbf{x}}_{\mathrm{WLS}}\right) \mathbf{R}^{-1} \mathbf{G}_{\mathbf{x}} \left( \widehat{\mathbf{x}}_{\mathrm{WLS}} \right)\right)^{-1}.
\end{equation}
This approximation is useful, but it should be interpreted as a local approximation. In strongly nonlinear problems, the quality of the covariance approximation depends on the accuracy of the linearisation and on the quality of the final estimate.

% ============================================
\subsection{Dynamic sensor fusion}
% ============================================

The preceding formulation treats $\mathbf{x}$ as static. In dynamic sensor fusion, the state evolves over time. A general state-space model consists of a process model and a measurement model. The process model is
\begin{equation}
\mathbf{x}_k
=
\mathbf{f}_{k-1}
\left(
\mathbf{x}_{k-1}
\right)
+
\mathbf{q}_{k-1},
\end{equation}
where $\mathbf{q}_{k-1}$ is process noise with covariance $\mathbf{Q}_{k-1}$. The measurement model is
\begin{equation}
\mathbf{y}_k
=
\mathbf{g}_k
\left(
\mathbf{x}_k
\right)
+
\mathbf{r}_k,
\end{equation}
where $\mathbf{r}_k$ is measurement noise with covariance $\mathbf{R}_k$.

Dynamic sensor fusion combines prediction and correction. The process model predicts how the state evolves, while the measurement model corrects that prediction using sensor data. In a batch least-squares formulation over a trajectory $\mathbf{x}_{0:K}$, one may write
\begin{align}
J(\mathbf{x}_{0:K})
=&
\left\|
\mathbf{x}_0
-
\mathbf{m}_0
\right\|_{\mathbf{P}_0^{-1}}^2
+
\sum_{k=1}^{K}
\left\|
\mathbf{x}_k
-
\mathbf{f}_{k-1}
\left(
\mathbf{x}_{k-1}
\right)
\right\|_{\mathbf{Q}_{k-1}^{-1}}^2 \nonumber \\
&+
\sum_{k=1}^{K}
\left\|
\mathbf{y}_k
-
\mathbf{g}_k
\left(
\mathbf{x}_k
\right)
\right\|_{\mathbf{R}_k^{-1}}^2.
\end{align}
The first term encodes prior information about the initial state. The second term enforces consistency with the dynamics. The third term enforces consistency with the sensor measurements. In this view, filtering and smoothing are extensions of regularised least squares to time-varying systems.



% GAB: I reduced it a bit as it was too long and a bit redundant. 
\subsection{Concluding perspective}

Two ideas from the preceding derivations survive into every method in this \deleted{survey}\added{book}. The first is that a measurement means nothing without its model and its uncertainty. What a sensor reports does not by itself determine how much it should move the estimate; that is fixed by what the sensor observes, how sensitively it responds to the state, how noisy it is, and whether it is biased. This is why covariance matrices are not bookkeeping: they encode how much each measurement and each direction in measurement space is to be believed.

The second is that independence is what makes information additive. Two independent measurements contribute genuinely separate evidence; two correlated ones share part of theirs, and treating them as independent counts the shared part twice. The estimate then reports a confidence it has not earned. This is the characteristic failure mode of sensor fusion, and it is worse than plain inaccuracy, because a fused estimate that is confidently wrong will be acted upon. The remedy is not fewer sensors but honest models: an ignored bias, a miscalibrated covariance, or an unmodelled dependency degrades the fused estimate below what a single well-understood sensor would have delivered.

Sequential estimation adds the operational reading. Past measurements need not be stored, because their effect is summarised by the current estimate and its covariance; a new measurement is then weighed by how surprising it is, through the innovation, and by how reliable it is, through the gain. This prediction--correction structure is the bridge from static least squares to the recursive filters of the next subsection.
This is the essence of modern sensor fusion: \emph{not merely combining sensor outputs, but combining information with an explicit account of uncertainty.}



%OLD VERSION
%% ============================================
%\subsection{Concluding perspective}
%% ============================================
%
%The mathematical foundation of sensor fusion is the idea that measurements are not used as isolated numerical readings, but as uncertain pieces of information about an underlying quantity of interest. Thus, sensor fusion is fundamentally an inverse problem under uncertainty: the measurements are observed, the state is hidden, and the task is to infer which state best explains the data while respecting the reliability of each source. The key point is that a sensor measurement is meaningful only together with its model and uncertainty description. A number reported by a sensor does not, by itself, say how strongly it should influence the estimate. Its influence depends on what physical quantity it observes, how sensitive it is to the unknown state, how noisy it is, whether its error is biased or zero mean, and whether it shares uncertainty with other measurements. This is why covariance matrices play such a central role. They are not merely technical parameters; they encode the confidence assigned to different measurements and directions in the measurement space.
%
%%
%
%%
%
%This viewpoint also explains why independence and correlation are critical. If two measurements are independent, their information can be combined as genuinely new evidence. If they are correlated, part of their information is shared. Treating correlated measurements as independent leads to double counting. The result is often an estimate that appears more accurate than it really is. In sensor fusion, this is one of the most important failure modes: the estimate may look statistically confident while being physically unreliable.
%
%Prior information has a similar interpretation. A prior is not an artificial mathematical addition; it represents information about the state that is available before the current measurements are processed. It may come from a previous estimate, a physical constraint, a map, a model, or domain knowledge. Regularisation incorporates this information by balancing consistency with the measurements against consistency with what is already known. The strength of the prior is determined by its uncertainty. A strong prior should dominate only when it is trustworthy; a weak prior should be easily corrected by informative measurements.
%
%Sequential estimation shows that fusion is also an information-processing mechanism. Past measurements need not always be stored explicitly. Their effect can be summarised by the current estimate and its covariance. A new measurement is then used to correct the previous estimate according to how surprising it is and how reliable it is. The innovation measures disagreement between prediction and observation, while the gain determines how much that disagreement should change the estimate. This prediction-correction interpretation is the conceptual bridge between static least squares and recursive filters such as the Kalman filter.
%
%%
%
%The main lesson is that sensor fusion is not about blindly combining more measurements. More data are useful only when their models, uncertainties, and dependencies are represented correctly. A poorly calibrated sensor, an incorrect covariance matrix, an ignored bias, or an unmodelled correlation can make the fused estimate worse rather than better. Conversely, when the models and uncertainty descriptions are credible, fusion can exploit complementarity: one sensor may compensate for another sensor's weakness, reduce ambiguity, improve robustness, and quantify the remaining uncertainty.
%
%The mathematical foundation of sensor fusion is therefore the careful conversion of heterogeneous measurements into information about the unknown state. Each measurement contributes according to what it observes, how accurately it observes it, and how independent it is from the other information sources. 
%This is the essence of modern sensor fusion: \emph{not merely combining sensor outputs, but combining information with an explicit account of uncertainty}.
%
%% ============================================
%%



\section{Bayesian Filtering}
\label{sec:bayes-filtering}
%
% ============================================

%Bayesian filtering provides a principled framework for estimating hidden states from noisy and incomplete measurements~\cite{sarkka2013_bayesian}. In autonomous driving, the hidden state may represent the ego-vehicle pose and velocity, the state of a surrounding object, the occupancy of a grid cell, or the motion of a tracked agent. The general Bayesian filtering process alternates between a prediction step, which propagates the state estimate using a motion model, and an update step, which incorporates new sensor measurements.


Bayesian filtering provides a probabilistic framework for recursively estimating a time-varying hidden state from noisy, incomplete, and possibly heterogeneous measurements~\cite{sarkka2013_bayesian}. For example, in autonomous driving, the state may represent the ego-vehicle pose, velocity, and inertial-sensor biases; the position, velocity, heading, and extent of a surrounding object; or the occupancy state of a spatial cell. Unlike static estimation, Bayesian filtering explicitly combines a dynamical model with measurements arriving over time.

Let $\mathbf{x}_k\in\mathbb{R}^{n_x}$ denote the hidden state at time $k$, $\mathbf{u}_{k-1}$ a known control or proprioceptive input, and $\mathbf{y}_k\in\mathbb{R}^{n_y}$ the measurements available at that time. A general state-space model is
\begin{subequations}
\begin{align}
\mathbf{x}_k &= \mathbf{f}_{k-1}(\mathbf{x}_{k-1},\mathbf{u}_{k-1})+\mathbf{q}_{k-1}, \label{eq:bayes_process_model}\\
\mathbf{y}_k &= \mathbf{g}_k(\mathbf{x}_k)+\mathbf{r}_k,
\label{eq:bayes_measurement_model}
\end{align}
\end{subequations}
where $\mathbf{q}_{k-1}$ and $\mathbf{r}_k$ denote process and measurement uncertainty, respectively. The process model describes how the state evolves, whereas the measurement model describes how the current state is observed by the sensors. The process noise accounts for manoeuvres, wheel slip, discretisation errors, and other imperfections in the motion model. The measurement noise represents sensor noise, calibration uncertainty, finite resolution, and modelling errors in the observation process. The filtering objective is to compute the posterior density
\begin{equation}
p(\mathbf{x}_k\mid\mathbf{y}_{1:k},\mathbf{u}_{0:k-1}),
\label{eq:filtering_posterior}
\end{equation}
where $\mathbf{y}_{1:k}\triangleq \{\mathbf{y}_1,\ldots,\mathbf{y}_k\}$. Under the usual first-order Markov assumption,
\begin{equation}
p(\mathbf{x}_k\mid\mathbf{x}_{0:k-1},\mathbf{u}_{0:k-1}) = p(\mathbf{x}_k\mid\mathbf{x}_{k-1},\mathbf{u}_{k-1}),
\end{equation}
and the conditional-independence assumption
\begin{equation}
p(\mathbf{y}_k\mid\mathbf{x}_{0:k},\mathbf{y}_{1:k-1}) = p(\mathbf{y}_k\mid\mathbf{x}_k),
\end{equation}
the posterior can be propagated recursively through prediction and
measurement-update steps. The \emph{prediction step} propagates the previous posterior through the state-transition model given by
\begin{align}
p&(\mathbf{x}_k\mid\mathbf{y}_{1:k-1}) \label{eq:bayes_prediction} \\
&= \int p(\mathbf{x}_k\mid\mathbf{x}_{k-1},\mathbf{u}_{k-1})
p(\mathbf{x}_{k-1}\mid\mathbf{y}_{1:k-1})
\,\mathrm{d}\mathbf{x}_{k-1}. \nonumber
\end{align}
This density is the prior belief about the current state before the new
measurement is processed. The \emph{update step} then applies Bayes'
rule:
\begin{equation}
p(\mathbf{x}_k\mid\mathbf{y}_{1:k}) = \frac{
p(\mathbf{y}_k\mid\mathbf{x}_k)
p(\mathbf{x}_k\mid\mathbf{y}_{1:k-1})
}{
p(\mathbf{y}_k\mid\mathbf{y}_{1:k-1})
},
\label{eq:bayes_update}
\end{equation}
where
\begin{equation}
p(\mathbf{y}_k\mid\mathbf{y}_{1:k-1}) = \int
p(\mathbf{y}_k\mid\mathbf{x}_k)
p(\mathbf{x}_k\mid\mathbf{y}_{1:k-1})
\,\mathrm{d}\mathbf{x}_k
\end{equation}
is the predictive likelihood, or evidence, of the new measurement. Hence, the posterior is obtained by weighting the predicted state distribution according to how well each candidate state explains the new sensor observation. This recursion provides the general mathematical foundation of Bayesian filtering. Different filters arise from different representations and approximations of the posterior distribution. 


%Kalman filtering and its variants are among the most widely used Bayesian methods in sensor fusion. The standard Kalman filter is appropriate for linear systems with Gaussian noise~\cite{kalman_1960}, while the extended Kalman filter linearises nonlinear models around the current estimate. The unscented Kalman filter uses deterministic sigma points to approximate nonlinear transformations~\cite{julier2004_ukf}, and particle filters represent the posterior distribution using weighted samples~\cite{gordon1993_pf,arulampalam2002_pf_tutorial}. These methods are widely used for GNSS--IMU fusion, vehicle state estimation, radar tracking, LiDAR-based object tracking, and camera--radar association.

When the process and measurement models are linear,
\begin{subequations}
\begin{align}
\mathbf{x}_k &= \mathbf{F}_{k-1}\mathbf{x}_{k-1} +\mathbf{B}_{k-1}\mathbf{u}_{k-1} + \mathbf{q}_{k-1}, \label{eq:kf_linear_process}\\
\mathbf{y}_k &= \mathbf{H}_k\mathbf{x}_k+\mathbf{r}_k,
\label{eq:kf_linear_measurement}
\end{align}    
\end{subequations}
and the noises are Gaussian, i.e, $\mathbf{q}_{k-1}\sim \mathcal{N}(\mathbf{0},\mathbf{Q}_{k-1})$ and $\mathbf{r}_k\sim \mathcal{N}(\mathbf{0},\mathbf{R}_k)$, the posterior remains Gaussian. It can therefore be represented exactly by its conditional mean and covariance, that is,
\begin{equation}
p(\mathbf{x}_k\mid\mathbf{y}_{1:k}) = \mathcal{N} \bigl(\widehat{\mathbf{x}}_{k|k}, \mathbf{P}_{k|k} \bigr).
\end{equation}
The resulting recursion is the Kalman filter~\cite{kalman_1960}. Its
prediction equations are
\begin{align}
\widehat{\mathbf{x}}_{k|k-1}
&=
\mathbf{F}_{k-1}\widehat{\mathbf{x}}_{k-1|k-1}
+
\mathbf{B}_{k-1}\mathbf{u}_{k-1},
\label{eq:kf_state_prediction}\\
\mathbf{P}_{k|k-1}
&=
\mathbf{F}_{k-1}\mathbf{P}_{k-1|k-1}
\mathbf{F}_{k-1}^{\mathsf T}
+
\mathbf{Q}_{k-1}.
\label{eq:kf_cov_prediction}
\end{align}
The new measurement is summarised by the innovation
\begin{equation}
\boldsymbol{\nu}_k
=
\mathbf{y}_k-
\mathbf{H}_k\widehat{\mathbf{x}}_{k|k-1},
\label{eq:kf_innovation}
\end{equation}
whose covariance is
\begin{equation}
\mathbf{S}_k
=
\mathbf{H}_k\mathbf{P}_{k|k-1}\mathbf{H}_k^{\mathsf T}
+
\mathbf{R}_k.
\label{eq:kf_innovation_covariance}
\end{equation}
The Kalman gain,
\begin{equation}
\mathbf{K}_k
=
\mathbf{P}_{k|k-1}\mathbf{H}_k^{\mathsf T}
\mathbf{S}_k^{-1},
\label{eq:kf_gain}
\end{equation}
balances uncertainty in the prediction against uncertainty in the new
measurement. The posterior mean and covariance are then updated as
\begin{align}
\widehat{\mathbf{x}}_{k|k} &= \widehat{\mathbf{x}}_{k|k-1} + \mathbf{K}_k\boldsymbol{\nu}_k, \label{eq:kf_state_update}\\
\mathbf{P}_{k|k} &= (\mathbf{I}-\mathbf{K}_k\mathbf{H}_k) \mathbf{P}_{k|k-1}. \label{eq:kf_cov_update}
\end{align} 
These equations have the same information-fusion interpretation as the sequential least-squares recursion developed in Section~\ref{subsubsec:sequential_least_squares}: the predicted state plays the role of a prior, the innovation measures disagreement between the prediction and the observation, and the gain determines how strongly the new information modifies the estimate.

When the models in \eqref{eq:bayes_process_model}--\eqref{eq:bayes_measurement_model} are nonlinear, the integrals in \eqref{eq:bayes_prediction} and
\eqref{eq:bayes_update} generally cannot be evaluated in closed form.
The extended Kalman filter (EKF) retains a Gaussian approximation and
linearises the models around the current estimate. Defining
\begin{subequations}
\begin{align}
\mathbf{F}_{k-1}
&=
\left.
\frac{\partial\mathbf{f}_{k-1}}
{\partial\mathbf{x}}
\right|_{\widehat{\mathbf{x}}_{k-1|k-1},
\mathbf{u}_{k-1}},
\\
\mathbf{H}_k
&=
\left.
\frac{\partial\mathbf{g}_k}
{\partial\mathbf{x}}
\right|_{\widehat{\mathbf{x}}_{k|k-1}},
\end{align}
\end{subequations}
the EKF applies the Kalman covariance and gain equations using these
local Jacobians, while evaluating the state and measurement predictions
through the original nonlinear functions. Its performance therefore
depends on the accuracy of the local linearisation; strong
nonlinearities or a poor initial estimate may lead to biased estimates,
inconsistent covariances, or divergence.

The unscented Kalman filter (UKF) avoids explicit
linearisation~\cite{julier2004_ukf}. It selects a deterministic set of
sigma points whose weighted mean and covariance match the current
Gaussian approximation, propagates these points through the nonlinear
models, and reconstructs the predicted moments from the transformed
points. The UKF can capture nonlinear transformations more accurately
than first-order linearisation, although it still represents the
posterior by a single Gaussian distribution. Particle filters provide a more general approximation for nonlinear and non-Gaussian problems~\cite{gordon1993_pf,arulampalam2002_pf_tutorial}. They represent the posterior by a weighted empirical distribution
\begin{equation}
p(\mathbf{x}_k\mid\mathbf{y}_{1:k})
\approx
\sum_{\ell=1}^{N_p}
w_k^{(\ell)}
\delta\!\left(
\mathbf{x}_k-\mathbf{x}_k^{(\ell)}
\right),
\label{eq:particle_posterior}
\end{equation}
where $\mathbf{x}_k^{(\ell)}$ is the $\ell$th particle,
$w_k^{(\ell)}\geq 0$, and
$\sum_{\ell=1}^{N_p}w_k^{(\ell)}=1$. Each particle is propagated
through the process model, and its weight is updated according to the
measurement likelihood:
\begin{equation}
\widetilde{w}_k^{(\ell)}
\propto
w_{k-1}^{(\ell)}
p\!\left(
\mathbf{y}_k\mid\mathbf{x}_k^{(\ell)}
\right).
\end{equation}
After normalisation, resampling is commonly used to remove particles
with negligible weight and replicate particles in regions of high
posterior probability. Particle filters can represent skewed,
multimodal, and non-Gaussian beliefs, but their computational cost
increases rapidly with the state dimension and the number of particles.


%For ego-localisation, Bayesian filters can combine high-rate inertial measurements with lower-rate GNSS updates, wheel odometry, visual odometry, or LiDAR odometry. The prediction step captures short-term motion, while external measurements correct accumulated drift. For object tracking, filters estimate object position, velocity, heading, and uncertainty over time. In this setting, the fusion problem is closely coupled with data association, because detections from different sensors and different time steps must be assigned to the correct tracks.

For ego-localisation, the state can, for example, contain position,
velocity, orientation, and IMU biases. High-rate IMU and wheel-odometry measurements propagate the state between external corrections, while GNSS, visual odometry, LiDAR
odometry, lane features, or map-matching constraints reduce accumulated drift. The covariance is propagated together with the state and quantifies how uncertainty grows during dead reckoning and contracts when informative external measurements become available.

For object tracking, a simple constant-velocity state may be characterized only by the position and velocity, although practical vehicle models may additionally include heading, turn rate, acceleration, dimensions, and object-existence variables. Camera, radar, and LiDAR detections then provide partial observations of this state. Radar may directly observe range and radial velocity, LiDAR may provide accurate position and extent, and cameras may provide bearing, class, and appearance information. The Bayesian update fuses these complementary observations according to their likelihoods and uncertainties.

In multi-object tracking, state estimation is inseparable from
\emph{data association}. Before a measurement can update a track, the
system must determine whether it originated from that track, from
another object, from a newly appearing object, or from clutter. A common
first step is innovation-based gating:
\begin{equation}
d_{k,j}^{2}
=
\boldsymbol{\nu}_{k,j}^{\mathsf T}
\mathbf{S}_{k,j}^{-1}
\boldsymbol{\nu}_{k,j}
\leq \gamma,
\label{eq:mahalanobis_gate}
\end{equation}
where $\boldsymbol{\nu}_{k,j}$ is the innovation obtained by assigning
measurement $j$ to the track and $\gamma$ is selected from a
$\chi^2$ distribution according to the measurement dimension and the
desired gate probability. Measurements outside the gate are considered
statistically incompatible with the predicted track. When several
measurements remain plausible, nearest-neighbour association,
probabilistic data association, joint probabilistic data association,
or multi-hypothesis tracking may be used to retain different degrees of
association uncertainty.



%A practical caveat concerns filter consistency: the reported covariance is trustworthy only if the process- and measurement-noise models are credible, and standard diagnostics such as the normalised innovation squared (NIS) test can be monitored online to detect model mismatch \cite{barshalom2001_estimation}. This diagnostic viewpoint recurs throughout the paper, from the innovation-based maritime monitoring of Chapter~\ref{sec:applications} to the hybrid classical--learning consistency checks of Chapter~\ref{sec:learning}. % [July 2026 pass] Added a consistency/diagnostics remark and citations; this subsection previously had no references at all.

A practical caveat concerns \emph{filter consistency}. The covariance
reported by a Bayesian filter is meaningful only when the dynamical
model, measurement model, process-noise covariance, measurement-noise
covariance, and assumed dependencies are credible. If uncertainty is
underestimated, the filter may reject valid measurements, react too
weakly to new information, or report an unrealistically precise state.
Conversely, excessive covariance produces unnecessarily noisy or
sluggish estimates. A standard online diagnostic is the normalised
innovation squared (NIS),
\begin{equation}
d^2_k
=
\boldsymbol{\nu}_k^{\mathsf T}
\mathbf{S}_k^{-1}
\boldsymbol{\nu}_k.
\label{eq:nis}
\end{equation}
Under a correctly specified linear--Gaussian model with an
$n_y$-dimensional measurement, $d^2_k$ follows a
$\chi^2_{n_y}$ distribution. Persistent violations of the corresponding
confidence bounds can indicate incorrect noise covariances, bias,
calibration drift, poor data association, sensor degradation, or
unmodelled dynamics~\cite{barshalom2001_estimation}. Innovation
whiteness, mean, and temporal correlation should also be monitored,
because a sequence may have approximately correct marginal variance
while still containing systematic or coloured modelling errors.

Bayesian filtering should therefore be understood not merely as a collection of recursive equations, but as a mechanism for propagating and revising beliefs about a hidden dynamic state. The prediction step encodes what the motion model implies before the new observation is seen; the update step determines how the observation changes that belief; and the posterior uncertainty records what remains unknown after both sources of information have been combined. This prediction--innovation--correction structure underlies localisation, tracking, navigation, occupancy estimation, and many hybrid classical--learning fusion architectures considered throughout this book.

Let us conclude this subsection with some examples.


\begin{example}
\label{ex:ekf-radar-cutin}
Consider an ego vehicle travelling at $27$\,m/s and a target vehicle travelling
at $22$\,m/s in the adjacent lane, initially $60$\,m ahead. All quantities are
expressed in the ego-vehicle frame, with $x$ forward and $y$ to the left and the
origin at the front-bumper radar. The state
\begin{equation}
\mathbf{x}_k=
\begin{bmatrix}
p_{x,k} & p_{y,k} & v_{x,k} & v_{y,k}
\end{bmatrix}^{\mathsf T}
\end{equation}
collects the relative position and velocity of the target, and is propagated by
the constant-velocity model of the form
\eqref{eq:kf_linear_process}, with
\begin{equation}
\mathbf{F}=
\begin{bmatrix}
1&0&T&0\\
0&1&0&T\\
0&0&1&0\\
0&0&0&1
\end{bmatrix},
\qquad T=0.05\ \mathrm{s},
\end{equation}
and a discretised white-noise-acceleration process covariance
\begin{equation}
\mathbf{Q}=\sigma_a^{2}
\begin{bmatrix}
\tfrac{T^{4}}{4}&0&\tfrac{T^{3}}{2}&0\\[2pt]
0&\tfrac{T^{4}}{4}&0&\tfrac{T^{3}}{2}\\[2pt]
\tfrac{T^{3}}{2}&0&T^{2}&0\\[2pt]
0&\tfrac{T^{3}}{2}&0&T^{2}
\end{bmatrix},
\qquad \sigma_a=0.8\ \mathrm{m/s^{2}},
\end{equation}
in which $\sigma_a$ expresses how much unmodelled acceleration the target is
permitted between updates. 
Let us assume the availability of a radar, which observes
\begin{equation}
\mathbf{g}(\mathbf{x}_k)=
\begin{bmatrix}
\sqrt{p_x^{2}+p_y^{2}}\\[4pt]
\arctan\!\left(p_y/p_x\right)\\[4pt]
\dfrac{p_xv_x+p_yv_y}{\sqrt{p_x^{2}+p_y^{2}}}
\end{bmatrix},
\qquad
\mathbf{R}=\operatorname{diag}\!\left(\sigma_r^{2},\sigma_\theta^{2},\sigma_{\dot r}^{2}\right),
\label{eq:ekf-radar-meas}
\end{equation}
with $\sigma_r=0.30$\,m, $\sigma_\theta=0.5^{\circ}$, and
$\sigma_{\dot r}=0.10$\,m/s, values representative of an automotive
$77$\,GHz radar. Writing $d=\sqrt{p_x^{2}+p_y^{2}}$, the Jacobian required by
the EKF is
\begin{equation}
\mathbf{H}_k=
\begin{bmatrix}
\dfrac{p_x}{d} & \dfrac{p_y}{d} & 0 & 0\\[8pt]
-\dfrac{p_y}{d^{2}} & \dfrac{p_x}{d^{2}} & 0 & 0\\[8pt]
\dfrac{p_y(v_xp_y-v_yp_x)}{d^{3}} &
\dfrac{p_x(v_yp_x-v_xp_y)}{d^{3}} &
\dfrac{p_x}{d} & \dfrac{p_y}{d}
\end{bmatrix},
\label{eq:ekf-radar-jacobian}
\end{equation}
evaluated at $\widehat{\mathbf{x}}_{k|k-1}$. Three practical points are worth
noting, because each of them is a common source of divergence. First,
\eqref{eq:ekf-radar-jacobian} is singular at $d=0$, so a target crossing the
sensor origin must be handled separately. Second, the azimuth component of the
innovation \eqref{eq:kf_innovation} must be wrapped to $(-\pi,\pi]$ before it is
used, since the difference of two angles is not the angle of the difference.
Third, the track is initialised directly from the first detection,
$\widehat{p}_x=r\cos\theta$ and $\widehat{p}_y=r\sin\theta$, with the velocity
components set to zero and given a deliberately large prior variance; the
azimuth channel alone cannot resolve the velocity, and pretending otherwise
produces an overconfident initial covariance.

The scenario contains one event the model does not: between $t=5$\,s and
$t=8$\,s the target changes lane into the path of the ego vehicle, moving
laterally from $3.5$\,m to $0$\,m. Fig.~\ref{fig:ekf-radar-cutin} shows the
outcome over a $10$\,s run at $20$\,Hz. In the upper panel the raw detections,
converted to Cartesian coordinates for display, scatter visibly about the true
trajectory, chiefly in the lateral direction: at $60$\,m an azimuth error of
$0.5^{\circ}$ is already a lateral error of $0.5$\,m, which is why a bearing-only
lateral estimate degrades with range while the range estimate does not. Over the
run the EKF reduces the position RMSE from $0.42$\,m to $0.16$\,m, a reduction of
$62$\,\%, and it does so while also estimating the two velocity components, which
the radar never observes directly.

The lower panel shows why this example is placed after the discussion of filter
consistency. It plots the normalised innovation squared \eqref{eq:nis} against
the two-sided $95\,\%$ interval of the $\chi^{2}_{3}$ distribution. While the
constant-velocity assumption holds, the filter is statistically consistent: the
mean NIS is $2.57$ against an expected value of $3$, and $98.8\,\%$ of the
samples fall inside the interval. During the lane change the picture changes.
The mean NIS rises to $3.73$, only $90.2\,\%$ of the samples remain inside the
interval, and the peak reaches $13.5$. Nothing is wrong with the sensor; what
the diagnostic has detected is that the process model no longer describes the
target. The estimate degrades accordingly but gracefully, with a peak lateral
error of $0.40$\,m and a lateral velocity estimate that lags the true one by
roughly a factor of two at the middle of the manoeuvre, which is the
characteristic signature of a constant-velocity filter tracking a manoeuvring
object.

This is the practical lesson the example is meant to carry. The EKF is not
failing here, and enlarging $\mathbf{Q}$ until the NIS looks acceptable would
only trade the lag for a noisier estimate in the much longer intervals during
which the target is not manoeuvring. The structural answer is to stop insisting
on a single motion model, which is precisely the argument for the interacting
multiple-model filters mentioned earlier in this subsection, and for treating
the innovation sequence as a monitored signal in its own right rather than as an
intermediate quantity.
\end{example}

\begin{figure}[!t]
\centering
\resizebox{0.99\columnwidth}{!}{
\begin{tikzpicture}
\begin{groupplot}[
    group style={group size=1 by 2, vertical sep=17mm},
    width=0.95\columnwidth,
    scale only axis,
    tick label style={font=\scriptsize},
    label style={font=\scriptsize},
    legend style={font=\scriptsize, draw=none, fill=white, fill opacity=0.85,
                  text opacity=1},
    grid=major,
    grid style={line width=0.2pt, draw=gray!25},
    axis line style={line width=0.5pt},
    tick align=outside,
    title style={font=\scriptsize, yshift=0.5mm},
]

%% ---------------- panel (a): bird's-eye view ----------------
\nextgroupplot[
    height=3.3cm,
    xlabel={longitudinal range $p_x$ [m]},
    ylabel={lateral offset $p_y$ [m]},
    xmin=8, xmax=63,
    ymin=-1.4, ymax=6.4,
    ytick={0,1,2,3,4,5},
    legend pos=north west,
    legend cell align=left,
    legend columns=3,
    title={(a) bird's-eye view in the ego-vehicle frame},
]
% lane divider the target crosses (ego lane centred on y = 0, adjacent lane on y = 3.5)
\addplot[draw=gray!55, dashed, line width=0.4pt, forget plot]
    coordinates {(8,1.75) (63,1.75)};

\addplot[only marks, mark=*, mark size=0.5pt, draw=MeasGrey, fill=MeasGrey]
    table[x=px_meas, y=py_meas, col sep=comma] {data/ch03/ekf_radar_cutin.csv};
\addlegendentry{radar detections}

\addplot[TruthBlack, line width=1.0pt]
    table[x=px_true, y=py_true, col sep=comma] {data/ch03/ekf_radar_cutin.csv};
\addlegendentry{ground truth}

\addplot[EkfBlue, line width=0.9pt, densely dashed]
    table[x=px_est, y=py_est, col sep=comma] {data/ch03/ekf_radar_cutin.csv};
\addlegendentry{EKF}

% the target closes on the ego vehicle, so time runs right to left
\draw[->, >=stealth, gray!70, line width=0.5pt]
    (axis cs:52,-0.75) -- (axis cs:20,-0.75);
\node[anchor=south, font=\scriptsize, text=black!60] at (axis cs:36,-0.9)
    {closing: $t=0$ at $p_x=60$\,m};

%% ---------------- panel (b): NIS consistency ----------------
\nextgroupplot[
    height=3.1cm,
    xlabel={time $t$ [s]},
    ylabel={NIS $d_k^{2}$},
    xmin=0, xmax=10,
    ymin=0, ymax=15,
    ytick={0,3,6,9,12,15},
    title={(b) innovation consistency against the $\chi^{2}_{3}$ interval},
]
% two-sided 95% interval of chi^2_3: [0.216, 9.348]  (drawn first, so the
% manoeuvre band below remains visible across the full height of the axis)
\addplot[draw=none, fill=BandGrey, forget plot]
    coordinates {(0,0.216) (10,0.216) (10,9.348) (0,9.348)} \closedcycle;
% lane-change window
\addplot[draw=none, fill=ManoeuvreOrange, fill opacity=0.45, forget plot]
    coordinates {(5,0) (8,0) (8,15) (5,15)} \closedcycle;
% expected value E[d^2] = 3
\addplot[draw=gray!70, line width=0.5pt, densely dotted, forget plot]
    coordinates {(0,3) (10,3)};

\addplot[EkfBlue, line width=0.5pt]
    table[x=t, y=nis, col sep=comma] {data/ch03/ekf_radar_cutin.csv};

\node[anchor=south, font=\scriptsize, text=black!70] at (axis cs:6.5,13.2) {lane change};
\node[anchor=west, font=\scriptsize, text=black!55] at (axis cs:0.2,10.4) {$95\,\%$ interval};

\end{groupplot}
\end{tikzpicture}
}
\caption{Extended Kalman filtering from automotive-radar
range, azimuth, and range-rate detections (Example~\ref{ex:ekf-radar-cutin}).
(a)~Raw detections converted to Cartesian coordinates, ground truth, and EKF
estimate. The target closes on the ego vehicle, so time runs from right to left;
the excursion at the right-hand edge is the initialisation transient, before the
velocity components have been observed. The detection scatter is dominated by
the azimuth channel and therefore grows with range. (b)~Normalised innovation
squared against the two-sided $95\,\%$ interval of $\chi^{2}_{3}$. The filter is
consistent while the constant-velocity assumption holds and becomes measurably
inconsistent during the lane change (shaded), which is the intended behaviour of
the diagnostic: it reports a model mismatch, not a sensor fault.}
\label{fig:ekf-radar-cutin}
\end{figure}

\begin{example}[Does the unscented transform help here?]
\label{ex:ukf-radar-cutin}
Example~\ref{ex:ekf-radar-cutin} is a natural place to ask whether the UKF
would do better, since it removes the first-order linearisation that the EKF
depends on. Running a UKF on the same measurement sequence, with the same
$\mathbf{F}$, $\mathbf{Q}$, and $\mathbf{R}$, answers the question. One
implementation detail is worth stating: the process model is linear, so the
prediction step remains the exact Kalman prediction and the unscented transform
is applied only to the measurement update. This is the usual choice when only
one of the two models is nonlinear, and it halves the number of sigma-point
propagations.

Over $200$ Monte Carlo runs the two filters are indistinguishable: the position
RMSE is $0.1775$\,m for the EKF and $0.1776$\,m for the UKF, the paired
difference averages $0.04\,\%$ of the RMSE, and the largest difference observed
on any single run is $9$\,mm. The reason is not that the measurement model is
mildly nonlinear in absolute terms, since the azimuth channel is decidedly not,
but that what governs the linearisation error is the curvature of
$\mathbf{g}$ \emph{over the support of the predicted covariance}. Here the
predicted lateral standard deviation settles around $0.11$\,m at ranges of tens
of metres, so the sigma points span an arc of roughly $0.2^{\circ}$, over which
$\arctan$ is affine to well below the noise floor. The EKF is not being
forgiven a poor approximation; there is no approximation error to speak of.

The comparison becomes interesting only when that support grows, which in
practice means sparse updates, long dead-reckoning intervals between external
corrections, or a poor initialisation. Table~\ref{tab:ekf-ukf-rate} varies the
radar update rate alone, leaving the geometry and all covariances unchanged.

\begin{table}[!t]
\centering
\small
\caption{EKF and UKF on the scenario of
Example~\ref{ex:ekf-radar-cutin} as the update rate is reduced, $100$ Monte
Carlo runs per row. The controlling quantity is $\sigma_y$, the predicted
lateral standard deviation entering the update, not the nonlinearity of
$\mathbf{g}$, which is the same in every row.}
\label{tab:ekf-ukf-rate}
\begin{tabular}{r c c c}
\hline
\textbf{Rate} & \textbf{$\sigma_y$ [m]} & \textbf{EKF RMSE [m]} & \textbf{UKF RMSE [m]}\\
\hline
$20$\,Hz  & $0.11$ & $0.098$ & $0.098$\\
$5$\,Hz   & $0.24$ & $0.188$ & $0.188$\\
$1$\,Hz   & $0.91$ & $0.335$ & $0.343$\\
$0.5$\,Hz & $2.56$ & $0.382$ & $0.508$\\
\hline
\end{tabular}
\end{table}

The UKF does not overtake the EKF anywhere in this problem, and at the lowest
rate it is clearly worse. That last row should be read with care rather than
as a general result: at $\sigma_y=2.56$\,m the sigma points straddle a genuinely
curved region of $\mathbf{g}$, the covariance update
$\mathbf{P}-\mathbf{K}\mathbf{P}_{zz}\mathbf{K}^{\mathsf T}$ has no Joseph-form
equivalent, and the outcome becomes sensitive to the sigma-point tuning. Two
practical remarks follow. First, that tuning is not innocuous: with the
frequently quoted $\alpha=10^{-3}$ the sigma points collapse to $0.002\sigma$
and the centre weight becomes $-8\times10^{5}$, so the UKF degenerates into a
central-difference EKF with poor conditioning; the values used here are
$\alpha=1$, $\beta=2$, $\kappa=1$, which keep all weights
positive~\cite{julier2004_ukf}. Second, the UKF costs $2n_x+1=9$ evaluations of
$\mathbf{g}$ per update against a single Jacobian. The conclusion for
automotive radar tracking is therefore narrow but useful: reach for the
unscented transform when the predicted uncertainty is large compared with the
scale on which the measurement model curves, and not merely because the
measurement model is nonlinear.
\end{example}







\section{Set-Membership and Zonotopic Estimation}
\label{sec:set-membership}

The probabilistic filters discussed above describe uncertainty through probability distributions, typically by propagating a mean and covariance or a set of weighted particles. In many sensor-fusion settings, however, only deterministic bounds are defensible. A sensor data sheet may provide a maximum error, a calibration procedure may certify a tolerance interval, or a safety monitor may need to guarantee that all states consistent with bounded disturbances remain inside an admissible region. Set-membership estimation addresses this case by propagating a set of all states compatible with the process model, measurements, and assumed error bounds, rather than a single most likely state \cite{alamo2005_zonotopes,rego2020_constrained_zonotopes}.

Zonotopes are a common representation for such bounded-error fusion because they are compact and closed under affine maps and Minkowski sums \cite{girard2005_zonotopes}. A zonotope in $\mathbb{R}^{d_x}$ can be written as
\begin{equation}
\mathcal{Z}=p\oplus H\mathbb{B}_{\infty}^{m}
=\left\{p+H\xi:\xi\in[-1,1]^m\right\},
\end{equation}
where $p$ is the centre, $H$ is the generator matrix, and $\mathbb{B}_{\infty}^{m}$ is the unit box. The centre can be interpreted as a nominal estimate, while the columns of $H$ describe the directions and magnitudes of admissible uncertainty. Compared with a full polytope, this centre--generator form is computationally attractive for embedded fusion because propagation mainly requires matrix products, column concatenation, and occasional order reduction.
The two algebraic operations used repeatedly by zonotopic estimators are exact in this representation. If $\mathcal{Z}_a=p_a\oplus H_a\mathbb{B}_{\infty}^{m_a}$ and $\mathcal{Z}_b=p_b\oplus H_b\mathbb{B}_{\infty}^{m_b}$, then
\begin{align}
A\mathcal{Z}_a+b
&=(Ap_a+b)\oplus AH_a\mathbb{B}_{\infty}^{m_a},\\
\mathcal{Z}_a\oplus\mathcal{Z}_b
&=(p_a+p_b)\oplus
\begin{bmatrix}H_a&H_b\end{bmatrix}
\mathbb{B}_{\infty}^{m_a+m_b}.
\end{align}
Thus, affine propagation changes the centre and generator directions, while adding independent uncertainty simply appends new generator columns.

Consider a dynamic sensor-fusion model
\begin{align}
\mathbf{x}_{k+1}&=\mathbf{f}_k(\mathbf{x}_k,\mathbf{u}_k)+\mathbf{w}_k,\\
\mathbf{y}_k&=\mathbf{h}_k(\mathbf{x}_k)+\mathbf{v}_k,
\end{align}
where $\mathbf{w}_k\in\mathcal{W}_k$ and $\mathbf{v}_k\in\mathcal{V}_k$ are unknown but bounded
\textcolor{olive}{, in contrast with the zero-mean stochastic $\mathbf{q}_k$ and $\mathbf{r}_k$ assumed earlier: the model is the same, only the description of the noise changes from a covariance to a set}. 

If the current state enclosure is $\mathcal{X}_k=p_k\oplus H_k\mathbb{B}_{\infty}^{m_k}$, the prediction step maps this set through the dynamics and adds the process-noise set. For nonlinear dynamics, a practical extended zonotopic estimator linearises the model around $p_k$ and appends a bounded Taylor-remainder zonotope:
\begin{equation}
\widetilde{\mathcal{X}}_{k+1}
\supseteq
\mathbf{f}_k(p_k,\mathbf{u}_k)
\oplus F_k H_k\mathbb{B}_{\infty}^{m_k}
\oplus \mathcal{Q}_k
\oplus \mathcal{E}^{f}_k,
\end{equation}
where $F_k=\partial \mathbf{f}_k/\partial \mathbf{x}|_{p_k}$ and $\mathcal{E}^{f}_k$ bounds the local linearisation error. This is analogous in form to the extended Kalman filter prediction, but the covariance update is replaced by a guaranteed set enclosure.

This prediction alone should be interpreted as a reachable-set propagation step, not yet as a complete estimator. It becomes zonotopic estimation only after the predicted enclosure is corrected by the measurements. To make this distinction explicit, let $\widetilde{\mathcal{X}}_{k+1}$ denote the predicted state set and let each scalar measurement channel define the measurement-consistent strip
\begin{equation}
\mathcal{S}_{i,k}
=
\left\{\mathbf{x}:\left|c_{i,k}^{\mathsf T}\mathbf{x}-\bar y_{i,k}\right|
\leq \bar\phi_{i,k}\right\},
\end{equation}
where $c_{i,k}^{\mathsf T}$ is the local measurement Jacobian row, $\bar y_{i,k}$ is the linearised measurement value, and $\bar\phi_{i,k}$ combines the sensor-noise bound with a nonlinear remainder bound. The exact set of states compatible with both the model prediction and the measurements is
\begin{equation}
\mathcal{F}_{k+1}
=
\widetilde{\mathcal{X}}_{k+1}
\cap
\bigcap_{i=1}^{d_y}\mathcal{S}_{i,k}.
\end{equation}
This expression is the set-membership analogue of the Bayesian update: the prediction supplies the prior enclosure, while each strip removes states inconsistent with a bounded-error sensor reading. The exact intersection $\mathcal{F}_{k+1}$ is generally a polytope rather than a zonotope. The filter therefore stores a corrected zonotope $\hat{\mathcal{X}}_{k+1}=\hat p\oplus\hat H\mathbb{B}_{\infty}^{\hat m}$ satisfying $\mathcal{F}_{k+1}\subseteq\hat{\mathcal{X}}_{k+1}$, so that the next prediction step can again use the efficient centre--generator calculus.

A standard multi-output strip correction gives such an outer enclosure directly from the predicted zonotope $\widetilde{\mathcal{X}}_{k+1}=\tilde p\oplus \tilde H\mathbb{B}_{\infty}^{\tilde m}$ \cite{le2013_zonotope_multioutput,narri2021_set_membership_shared_awareness}. For gain vectors $\lambda_{i,k}\in\mathbb{R}^{d_x}$, define
\begin{align}
\hat p
&=
\tilde p+
\sum_{i=1}^{d_y}
\lambda_{i,k}\left(\bar y_{i,k}-c_{i,k}^{\mathsf T}\tilde p\right),\\
H_{\mathrm{old}}
&=
\left(I-\sum_{i=1}^{d_y}\lambda_{i,k}c_{i,k}^{\mathsf T}\right)\tilde H,\\
H_{\mathrm{meas}}
&=
\begin{bmatrix}
\bar\phi_{1,k}\lambda_{1,k}&\cdots&\bar\phi_{d_y,k}\lambda_{d_y,k}
\end{bmatrix},\\
\hat H&=
\begin{bmatrix}
H_{\mathrm{old}}&H_{\mathrm{meas}}
\end{bmatrix}.
\end{align}
Here $\hat p$ is the centre of the corrected zonotope. Before any order reduction, $\hat H$ has $\hat m=\tilde m+d_y$ columns: $\tilde m$ inherited from the predicted enclosure and $d_y$ appended from the measurement strips. Each residual $\bar y_{i,k}-c_{i,k}^{\mathsf T}\tilde p$ measures the mismatch between the predicted centre and the $i$th measurement strip, and $\lambda_{i,k}$ determines how that residual shifts the centre toward the measurement-consistent region. The block $H_{\mathrm{old}}$ reshapes the previously predicted uncertainty directions after imposing all strip constraints, while $H_{\mathrm{meas}}$ appends one generator per measurement channel to preserve the full bounded-noise width. For a single strip this reduces to the familiar update
\begin{align}
\hat p&=\tilde p+\lambda_{i,k}\left(\bar y_{i,k}-c_{i,k}^{\mathsf T}\tilde p\right),\\
\hat H&=
\left[
\left(I-\lambda_{i,k}c_{i,k}^{\mathsf T}\right)\tilde H
\quad
\bar\phi_{i,k}\lambda_{i,k}
\right].
\end{align}
The gain is therefore not a stochastic Kalman gain, but a design variable chosen to reduce the size of the outer approximation while preserving containment.

\begin{figure}[t]

\centering
\resizebox{\columnwidth}{!}{%
\begin{tikzpicture}[scale=0.78,>=Stealth]
    \coordinate (p) at (0,0);
    \coordinate (v1) at ( 2.2,  2.9);
    \coordinate (v2) at ( 3.2,  1.3);
    \coordinate (v3) at ( 1.8, -1.5);
    \coordinate (v4) at (-2.2, -2.9);
    \coordinate (v5) at (-3.2, -1.3);
    \coordinate (v6) at (-1.8,  1.5);
    \def\yc{1.0}
    \def\phi{0.6}
    \coordinate (phat) at (0,\yc);
    \def\hx{3.2}
    \def\hy{\phi}

    \filldraw[fill=gray!16,draw=gray!65,thick]
      (v1)--(v2)--(v3)--(v4)--(v5)--(v6)--cycle;
    \path[fill=yellow!22,draw=yellow!55!black,thick]
      (-4,\yc+\phi) -- (4,\yc+\phi) -- (4,\yc-\phi) -- (-4,\yc-\phi) -- cycle;
    \begin{scope}
      \clip (v1)--(v2)--(v3)--(v4)--(v5)--(v6)--cycle;
      \fill[orange!70,opacity=0.82]
        (-4,\yc+\phi) -- (4,\yc+\phi) -- (4,\yc-\phi) -- (-4,\yc-\phi) -- cycle;
    \end{scope}
    \draw[gray!70,thick] (v1)--(v2)--(v3)--(v4)--(v5)--(v6)--cycle;
    \draw[yellow!55!black,thick] (-4,\yc+\phi)--(4,\yc+\phi);
    \draw[yellow!55!black,thick] (-4,\yc-\phi)--(4,\yc-\phi);
    \draw[dashed,orange!80!black,thick] (-4,\yc)--(4,\yc)
      node[right] {$c^\top x=y$};
    \draw[very thick,dashed,blue!75,fill=blue!8,fill opacity=0.42,draw opacity=1]
      ($ (phat)+(\hx,\hy)$) -- ($ (phat)+(\hx,-\hy)$) --
      ($ (phat)+(-\hx,-\hy)$) -- ($ (phat)+(-\hx,\hy)$) -- cycle;
    \fill (p) circle (1.7pt) node[below left] {$p$};
    \fill (phat) circle (1.7pt) node[above left] {$\hat p$};
    \node[draw,fill=gray!16,inner sep=2pt,anchor=west] at (-3.8,-3.45) {\scriptsize predicted $\widetilde{\mathcal X}_{k+1}$};
    \node[draw,fill=orange!70,inner sep=2pt,anchor=west] at (-0.7,-3.45) {\scriptsize exact intersection};
    \node[draw=blue!75,dashed,fill=blue!8,inner sep=2pt,anchor=west] at (2.25,-3.45) {\scriptsize kept zonotope};
\end{tikzpicture}}
\caption{Geometric interpretation of one strip correction. The grey set is the predicted zonotope, the yellow band is the bounded-error measurement strip, the orange region is the exact feasible intersection, and the dashed blue set is the zonotopic outer approximation propagated to the next step.}
\label{fig:zonotope_strip_correction}
\end{figure}
\begin{example}[Single-strip correction]
\label{ex:zonotope-strip}
Consider
\begin{equation}
p=
\begin{bmatrix}
0\\
0
\end{bmatrix},
\quad
H=
\begin{bmatrix}
2.0 & 0.7 & -0.5\\
0.7 & 1.4 & 0.8
\end{bmatrix}.
\end{equation}
With
\begin{equation}
c=
\begin{bmatrix}
0\\
1
\end{bmatrix},
\quad
y=1.0,\quad
\phi=0.6.
\end{equation}
The strip is $0.4\le x_2\le1.6$. With $\gamma=c$, the centre moves to $\hat p=[0\;1]^{\mathsf T}$, the corrected generator matrix is
\begin{equation}
\hat H=
\begin{bmatrix}
2.0 & 0.7 & -0.5 & 0\\
0 & 0 & 0 & 0.6
\end{bmatrix},
\end{equation}
and the retained zonotope encloses the exact intersection shown in Fig.~\ref{fig:zonotope_strip_correction}. \textcolor{olive}{The price of that guarantee is visible in the same figure: the retained set spans $x_1\in[-3.2,3.2]$, the full width of the prediction, because a zonotope cannot represent the slanted cut that the strip actually makes. Containment is exact; tightness is not.} Processing several sensor channels sequentially gives a corrected zonotope that still contains all measurement-consistent states, provided the assumed bounds are valid.
Processing several sensor channels sequentially gives a corrected zonotope that still contains all measurement-consistent states, provided the assumed bounds are valid.
\end{example}



\textcolor{olive}{Example~\ref{ex:zonotope-strip} shows such a correction mechanism from a numerical perspective.} The number of columns of the generator matrix is the number of independent uncertainty directions retained by the filter. Prediction appends process-noise and linearisation-remainder generators, and each strip correction appends a measurement-noise generator. Without compression, the matrix $H$ therefore grows at every time step, increasing memory use and making subsequent matrix products more expensive. Practical zonotopic estimators control this growth by order reduction. If
\begin{equation}
H=\left[H_{\mathrm{keep}}\quad H_{\mathrm{red}}\right],
\end{equation}
where $H_{\mathrm{keep}}$ contains the most important columns and $H_{\mathrm{red}}$ contains columns to be compressed, then a common enclosure-preserving reduction replaces $H_{\mathrm{red}}$ by an interval generator
\begin{equation}
D=\operatorname{diag}\!\left(\sum_j |(H_{\mathrm{red}})_{1j}|,\ldots,
\sum_j |(H_{\mathrm{red}})_{d_xj}|\right),
\end{equation}
and uses the reduced zonotope
\begin{equation}
\mathcal{Z}_{\mathrm{red}}
=
p\oplus
\left[H_{\mathrm{keep}}\quad D\right]\mathbb{B}_{\infty}.
\end{equation}
This replacement reduces many low-priority columns to only $d_x$ axis-aligned generators while preserving containment, because the discarded generators are enclosed by a box. The kept columns are typically selected according to their norm, their contribution to the zonotope size, or their alignment with dynamically important state directions. Thus the estimator trades a controlled amount of additional conservatism for bounded computational complexity.


% GAB: i shortened this as it vas too verbose.
\textcolor{olive}{A recent extension replaces the purely deterministic description with probabilistic zonotopes, in which a bounded component and a Gaussian component coexist in the same geometric object~\cite{modares2026reachability}. The motivation is that real uncertainty budgets are mixed: a calibration tolerance or a certified sensor limit is a hard bound, whereas thermal noise is not. The framework additionally separates disturbances already realised, which can be inferred from logged trajectories and used to tighten the data-consistent model set, from future aleatory ones, which must still be propagated to obtain high-probability reachable sets. The practical consequence for fusion is that data can be used not only to confirm that a conservative uncertainty model is safe, but to shrink it.}


%A recent extension of this line of work is the use of probabilistic zonotopes for data-driven reachability analysis and uncertainty refinement \cite{modares2026reachability}. While classical zonotopic estimators treat uncertainty mainly through deterministic bounded sets, probabilistic zonotopes combine bounded deterministic components with Gaussian stochastic components, thereby retaining the geometric advantages of zonotopes while also representing probabilistic uncertainty. This is particularly relevant for safety-oriented sensor fusion, where some uncertainty sources are best described by hard bounds, such as calibration tolerances or certified sensor error limits, whereas others are naturally stochastic, such as measurement noise or environmental variability. The framework in \cite{modares2026reachability} further distinguishes between realized disturbances, which are inferred from collected trajectories and used to tighten the data-consistent model set, and future aleatory disturbances, which must still be propagated to obtain high-probability reachable sets. This distinction is important for autonomous and robotic systems because data should not only validate whether a conservative prior uncertainty model is safe, but also refine it to reduce unnecessary conservatism. In this sense, probabilistic zonotopes provide a bridge between deterministic set-membership fusion and stochastic uncertainty propagation: they preserve containment-oriented reasoning while allowing data-driven refinement of the uncertainty model.


\begin{example}
\label{ex:zono-radar-cutin}
The machinery of this subsection can be applied to the scenario of
Example~\ref{ex:ekf-radar-cutin} with no change of model, which makes the
comparison between the two descriptions of uncertainty direct. The state and
the transition matrix $\mathbf{F}$ are the ones already defined there; what
changes is that the covariances $\mathbf{Q}$ and $\mathbf{R}$ are replaced by
hard bounds, and the filter propagates $\mathcal{X}_k=p_k\oplus H_k\mathbb{B}_{\infty}^{m_k}$
instead of a mean and a covariance.

Bounding the relative acceleration by $\bar a=2.5$\,m/s$^{2}$ makes the
process-noise set a two-generator zonotope, since a constant-velocity model
maps a bounded acceleration onto the two columns
$\bar a\,[\,T^{2}/2,\,0,\,T,\,0\,]^{\mathsf T}$ and
$\bar a\,[\,0,\,T^{2}/2,\,0,\,T\,]^{\mathsf T}$. The prediction step is then
exact: the dynamics are linear, so no Taylor-remainder zonotope is needed.

The measurement is where the construction is worth spelling out, because the
polar geometry that forced a Jacobian in Example~\ref{ex:ekf-radar-cutin} does
not force a linearisation here. Let $u=(\cos\theta,\sin\theta)$ and
$v=(-\sin\theta,\cos\theta)$ be built from the measured bearing. If the sensor
errors obey $|\Delta r|\le\bar r$ and $|\Delta\theta|\le\bar\theta$, then the
true position satisfies
\begin{align}
u^{\mathsf T}p &\in \left[(r-\bar r)\cos\bar\theta,\; r+\bar r\right],
\label{eq:zono-strip-radial}\\
v^{\mathsf T}p &\in \left[-(r+\bar r)\sin\bar\theta,\; (r+\bar r)\sin\bar\theta\right],
\label{eq:zono-strip-tangential}
\end{align}
which are two strips of exactly the form used above, with normals
$c_{1}=(u_x,u_y,0,0)$ and $c_{2}=(v_x,v_y,0,0)$. These are not approximations:
the annulus sector of admissible positions is contained in a rectangle in the
frame aligned with the bearing, and \eqref{eq:zono-strip-radial}--\eqref{eq:zono-strip-tangential}
are its two faces. The range-rate channel supplies a third strip with normal
$c_{3}=(0,0,u_x,u_y)$ and half-width $\bar{\dot r}+\bar\theta\,v_{\max}$, valid
whenever the relative speed is bounded by $v_{\max}$, since a bearing error
rotates the projection direction by at most $\bar\theta$. One further
observation ties the two viewpoints together: choosing the strip gain to
minimise the Frobenius radius of the corrected generator matrix gives the
closed form
\begin{equation}
\lambda^{\star}=\frac{\widetilde H\widetilde H^{\mathsf T}c}
{c^{\mathsf T}\widetilde H\widetilde H^{\mathsf T}c+\bar\phi^{2}},
\label{eq:zono-optimal-gain}
\end{equation}
which is the Kalman gain of Example~\ref{ex:ekf-radar-cutin} with
$\widetilde H\widetilde H^{\mathsf T}$ in the place of the covariance and
$\bar\phi^{2}$ in the place of the measurement noise.

Taking the sensor bounds at four standard deviations of the noise used in
Example~\ref{ex:ekf-radar-cutin}, that is $\bar r=1.2$\,m,
$\bar\theta=2^{\circ}$, $\bar{\dot r}=0.4$\,m/s, and running the estimator on
the same detection sequence gives Fig.~\ref{fig:zono-radar-cutin}. Every
assumption was checked rather than asserted: no realised sensor error exceeds
its bound over the $201$ steps, and the peak lateral acceleration of the lane
change is $2.33$\,m/s$^{2}$, inside $\bar a$. The result is the trade-off this
subsection exists to describe, in numbers. The zonotopic enclosure contains the
true state at \emph{every} step, including throughout the manoeuvre. The
$3\sigma$ envelope of the EKF does not: it is breached on $24$ of $201$ steps,
and the breaches are not spread evenly. During the constant-velocity segment it
is never violated; during the lane change it fails on $29.5\,\%$ of the steps.
This is the same model mismatch that the NIS reported in
Example~\ref{ex:ekf-radar-cutin}, seen from the other side, and it is worth
noting that the probabilistic envelope is only mildly wrong, the largest
excursion being $0.15$\,m.

The price is conservatism, and it is not small. The zonotopic lateral envelope
is between five and nine times wider than the EKF's $6\sigma$ band, narrowing
from $7.6$\,m at a range of $57$\,m to $1.4$\,m at $10$\,m. That profile is
itself informative: the width is dominated by the tangential strip
\eqref{eq:zono-strip-tangential}, whose half-width $(r+\bar r)\sin\bar\theta$
scales with range, so bounded-error tracking of a distant target is expensive
in exactly the way bearing-based lateral estimation is. Order reduction keeps
the generator count between $7$ and $20$. The practical reading is the one
already given above: use the enclosure where a guarantee is required and its
width is affordable, such as gating a learned detection or certifying a
collision-free corridor at short range, and keep the probabilistic estimate for
the point estimate that planning and control actually consume.
\end{example}

\begin{figure}[!t]
\centering
\begin{tikzpicture}
\begin{groupplot}[
    group style={group size=1 by 2, vertical sep=17mm},
    width=0.80\columnwidth,
    scale only axis,
    tick label style={font=\scriptsize},
    label style={font=\scriptsize},
    legend style={font=\scriptsize, draw=none, fill=white, fill opacity=0.85,
                  text opacity=1},
    grid=major,
    grid style={line width=0.2pt, draw=gray!25},
    axis line style={line width=0.5pt},
    tick align=outside,
    title style={font=\scriptsize, yshift=0.5mm},
    xmin=0, xmax=10,
]

%% ---------------- panel (a): the two envelopes, full scale ----------------
\nextgroupplot[
    height=3.1cm,
    xlabel={time $t$ [s]},
    ylabel={lateral offset $p_y$ [m]},
    ymin=-1.2, ymax=8.6,
    ytick={0,2,4,6,8},
    legend pos=north east,
    legend cell align=left,
    legend columns=2,
    title={(a) guaranteed enclosure against the $3\sigma$ envelope},
]
\addplot[draw=none, fill=ManoeuvreOrange, fill opacity=0.40, forget plot]
    coordinates {(5,-1.2) (8,-1.2) (8,8.6) (5,8.6)} \closedcycle;
\addplot[name path=zlo, draw=none, forget plot]
    table[x=t, y=zlo, col sep=comma] {data/ch03/zono_radar_cutin.csv};
\addplot[name path=zhi, draw=none, forget plot]
    table[x=t, y=zhi, col sep=comma] {data/ch03/zono_radar_cutin.csv};
\addplot[fill=ZonoGreen, fill opacity=0.18] fill between[of=zlo and zhi];
\addlegendentry{zonotope}
\addplot[name path=elo, draw=none, forget plot]
    table[x=t, y=ekf_lo, col sep=comma] {data/ch03/zono_radar_cutin.csv};
\addplot[name path=ehi, draw=none, forget plot]
    table[x=t, y=ekf_hi, col sep=comma] {data/ch03/zono_radar_cutin.csv};
\addplot[fill=EkfBlue, fill opacity=0.45] fill between[of=elo and ehi];
\addlegendentry{EKF $3\sigma$}
\addplot[TruthBlack, line width=0.8pt]
    table[x=t, y=py_true, col sep=comma] {data/ch03/zono_radar_cutin.csv};
\addlegendentry{ground truth}

%% ---------------- panel (b): the same, relative to the truth ----------------
\nextgroupplot[
    height=3.1cm,
    xlabel={time $t$ [s]},
    ylabel={error [m]},
    ymin=-0.65, ymax=0.65,
    ytick={-0.6,-0.3,0,0.3,0.6},
    title={(b) the same, relative to the truth},
]
\addplot[draw=none, fill=ManoeuvreOrange, fill opacity=0.40, forget plot]
    coordinates {(5,-0.65) (8,-0.65) (8,0.65) (5,0.65)} \closedcycle;
\addplot[name path=relo, draw=none, forget plot]
    table[x=t, y=re_lo, col sep=comma] {data/ch03/zono_radar_cutin.csv};
\addplot[name path=rehi, draw=none, forget plot]
    table[x=t, y=re_hi, col sep=comma] {data/ch03/zono_radar_cutin.csv};
\addplot[fill=EkfBlue, fill opacity=0.35, forget plot] fill between[of=relo and rehi];
\addplot[TruthBlack, line width=0.6pt, forget plot]
    coordinates {(0,0) (10,0)};
\addplot[EkfBlue, line width=0.6pt, forget plot]
    table[x=t, y=e_ekf, col sep=comma] {data/ch03/zono_radar_cutin.csv};
\node[anchor=west, font=\scriptsize, text=black!60] at (axis cs:0.25,-0.5)
    {band: EKF $3\sigma$; line: EKF error};

\end{groupplot}
\end{tikzpicture}
\caption{Bounded-error and probabilistic descriptions for Example~\ref{ex:zono-radar-cutin}, from the same detections.
(a)~The zonotopic enclosure contains the true lateral offset at every step; the
EKF $3\sigma$ envelope is far tighter but is not a guarantee. (b)~The same data
plotted relative to the truth: the $3\sigma$ band crosses zero on $29.5\,\%$ of
the steps inside the lane-change window (shaded) and on none of the steps while
the constant-velocity model holds. The conservatism of the enclosure is the
price of the guarantee, and it scales with range because the tangential strip
half-width is $(r+\bar r)\sin\bar\theta$.}
\label{fig:zono-radar-cutin}
\end{figure}

What Example~\ref{ex:zono-radar-cutin} does for a single tracked object
generalises, and three use cases recur in autonomous driving. Each makes the
bargain the example made explicit: an envelope several times wider than a
covariance, in exchange for a statement about what the system can guarantee
rather than what it expects. The first is ego-localisation under bounded
position, velocity, and map errors, where a GNSS--IMU--wheel-odometry localiser
propagates a zonotope through a vehicle-motion model and contracts it with
bounded GNSS, odometry, lane-boundary, or map-matching constraints, delivering
an ego-pose envelope rather than a pose. This matters most in urban GNSS, where
multipath and non-line-of-sight effects invalidate a purely Gaussian
description: set-valued shadow matching and mosaic zonotope formulations target
precisely that
regime~\cite{bhamidipati2022_zonotope_shadow_matching,neamati2023_mosaic_zonotope_shadow_matching}.

The second is conservative tracking and shared situational awareness. Radar,
LiDAR, camera, V2X, and roadside sensors each produce a bounded set for the same
object, and the fusion module intersects or outer-approximates them instead of
collapsing them to a point, which is attractive under occlusion, under ambiguous
association, or when an infrastructure sensor observes a region the ego vehicle
cannot see~\cite{narri2021_set_membership_shared_awareness,narri2025_set_based_awareness};
constrained and hybrid zonotopes extend this to inconsistent measurements and
confidence-weighted fusion in connected automated
vehicles~\cite{narri2026_hybrid_zonotope_awareness}.

The third is consistency monitoring around learned perception. A detection that
falls outside the reachable zonotope of an existing track, or whose induced strip
intersects the predicted set emptily or very conservatively, can be down-weighted,
retained as one of several hypotheses, or escalated to a conservative fallback.
Here the enclosure does not replace deep perception but bounds it, which is also
where set-membership methods meet learning-aware filters such as Gaussian-process
zonotopic Kalman filtering~\cite{li2022_gp_zkf}. All three inherit the caveat
the example already quantified: a guarantee is only as good as the assumed
disturbance, sensor, and calibration bounds, loose bounds buy conservatism while
overly tight ones may exclude the true state, and high-dimensional perception
states force aggressive generator reduction or a problem-specific choice of what
to enclose.


% ============================================
%
\section{Track-Level Fusion and Data Association}
\label{sec:track-level}
\label{subsec:track-fusion}
%
% ============================================

The statistical principles introduced in the preceding subsections also apply when the quantities being combined are not raw sensor measurements, but object-state estimates produced by separate tracking modules. This setting is known as \emph{track-level fusion}, or \emph{track-to-track fusion}. For example, radar, LiDAR, and camera subsystems may each detect and track the same vehicle independently, after which a higher-level fusion module associates the local tracks and combines their position, velocity, class, and uncertainty information \cite{liggins_handbook_2009,khaleghi_multisensor_2013,
barshalom1995multitarget,barshalom2011tracking}.

Track-level fusion differs from measurement-level fusion. In
measurement-level fusion, the original sensor observations, such as radar range and range rate, LiDAR centroids, or image bearings, are incorporated directly into a common estimator. In track-level fusion, each sensing subsystem first performs its own detection, association, filtering, and track management, and only the resulting track estimates are exchanged. This modular structure reduces communication and allows independently developed perception pipelines to be integrated. However, some information
is lost because the original measurements, local association decisions, and measurement histories are compressed into the reported tracks \cite{blackman1999design,barshalom2011tracking}.

% --------------------------------------------
\subsection{Local Track Representation and Alignment}
% --------------------------------------------

A local track produced by subsystem $i$ at time $k$ may be represented by
\begin{equation}
\mathcal{T}_{i,k}=\left(\widehat{\mathbf{x}}_{i,k},\mathbf{P}_{i,k},r_{i,k},c_{i,k},\tau_{i,k}\right),
\label{eq:local-track-description}
\end{equation}
where $\widehat{\mathbf{x}}_{i,k}$ and $\mathbf{P}_{i,k}$ are the local
state estimate and covariance, $r_{i,k}$ is an existence or track-quality
score, $c_{i,k}$ denotes class information, and $\tau_{i,k}$ is the
timestamp. Depending on the application, the track state may contain
position, velocity, acceleration, heading, turn rate, object dimensions,
or other kinematic and semantic quantities.

Before two tracks can be compared or fused, they must describe the same
state variables in a common coordinate frame and at a common reference
time. If subsystem $i$ reports its estimate at time $\tau_{i,k}$, it may
be propagated to the fusion time $t_k$ according to
\begin{equation}
\widehat{\mathbf{x}}_{i,k}^{\,\mathrm{ref}}=\mathbf{F}(t_k-\tau_{i,k})\widehat{\mathbf{x}}_{i,k},
\label{eq:track-time-alignment-state}
\end{equation}
and
\begin{equation}
\mathbf{P}_{i,k}^{\,\mathrm{ref}}=\mathbf{F}(t_k-\tau_{i,k})\mathbf{P}_{i,k}\mathbf{F}^{\mathsf T}(t_k-\tau_{i,k})+\mathbf{Q}(t_k-\tau_{i,k}),
\label{eq:track-time-alignment-cov}
\end{equation}
where $\mathbf{Q}(t_k-\tau_{i,k})$ accounts for the uncertainty accumulated
during the propagation interval. This covariance growth is essential in
asynchronous fusion: a delayed track should become less influential as its
age increases, unless new information supports its prediction
\cite{barshalom2011tracking}.

Spatial alignment is equally important. Tracks generated in a radar frame,
camera frame, LiDAR frame, ego-vehicle frame, roadside frame, or global map
frame must be transformed consistently. The associated covariance must also
be transformed through the corresponding Jacobian or linear coordinate
mapping. Otherwise, apparently precise tracks may be compared in
incompatible reference frames.

% --------------------------------------------
\subsection{Gating and Track Association}
% --------------------------------------------

Track-level fusion is meaningful only after establishing whether two local
tracks correspond to the same physical object. Let
$\widehat{\mathbf{x}}_{i,k}^{-}$ denote a predicted track from subsystem
$i$, and let $\widehat{\mathbf{x}}_{j,k}$ denote a candidate track from
subsystem $j$. If the two tracks use the same state representation, their
difference is
\begin{equation}
\boldsymbol{\delta}_{ij,k}=\widehat{\mathbf{x}}_{j,k}-\widehat{\mathbf{x}}_{i,k}^{-}.
\label{eq:track-difference}
\end{equation}

Under the preliminary assumption that the two errors are uncorrelated, the
covariance of the difference is
\begin{equation}
\mathbf{S}_{ij,k}=\mathbf{P}_{i,k}^{-}+\mathbf{P}_{j,k}.
\label{eq:track-difference-cov}
\end{equation}
The corresponding squared Mahalanobis distance is
\begin{equation}
d_{ij,k}^{2}=\boldsymbol{\delta}_{ij,k}^{\mathsf T}\mathbf{S}_{ij,k}^{-1}\boldsymbol{\delta}_{ij,k}.
\label{eq:track-mahalanobis}
\end{equation}
The pair is retained as a plausible association when
\begin{equation}
d_{ij,k}^{2}\leq\gamma,
\label{eq:track-gate}
\end{equation}
where $\gamma$ is selected from a $\chi^{2}$ distribution according to the
dimension of the compared state and the desired gate probability
\cite{blackman1999design,barshalom2011tracking}. Mahalanobis gating is
generally preferable to Euclidean gating because it accounts for both the
magnitude and the directional structure of the track uncertainty. For
example, two distant camera-based tracks may differ considerably in depth
but remain statistically compatible because their range uncertainty is
large.

Gating removes clearly implausible matches, but it does not resolve
competition when several assignments remain feasible. The principal
association methods differ in how much ambiguity they retain.

Global nearest-neighbour (GNN) association constructs a cost matrix from
the gated track pairs and selects one globally consistent assignment,
usually by solving a linear assignment problem with the Hungarian
algorithm. GNN is computationally attractive and suitable for real-time
systems, but it commits to a single hypothesis at each scan. In dense
traffic, clutter, partial occlusion, or crossing trajectories, an incorrect
early assignment may cause track fragmentation or identity switches
\cite{blackman1999design,barshalom2011tracking}.

Probabilistic data association (PDA) and joint probabilistic data
association (JPDA) retain uncertainty over the possible associations at
the current scan. Instead of updating a track using only one selected
candidate, JPDA forms a probability-weighted innovation from all candidates
inside its gate. If $\beta_{ij,k}$ denotes the probability that candidate
$j$ originated from track $i$, and $\beta_{i0,k}$ is the probability of a
missed detection, then
\begin{equation}
\overline{\boldsymbol{\nu}}_{i,k}=\sum_{j\in\mathcal{G}_{i,k}}\beta_{ij,k}\boldsymbol{\nu}_{ij,k},
\label{eq:jpda-innovation}
\end{equation}
with
\begin{equation}
\beta_{i0,k}+\sum_{j\in\mathcal{G}_{i,k}}\beta_{ij,k}=1,
\label{eq:jpda-probability}
\end{equation}
where $\mathcal{G}_{i,k}$ denotes the gated candidate set. JPDA is more
robust than nearest-neighbour association when several observations are
plausible, although neighbouring tracks may coalesce if they are repeatedly
updated using similar probability-weighted combinations
\cite{barshalom1995multitarget,blackman1999design,
barshalom2011tracking}.

Multi-hypothesis tracking (MHT) retains multiple association histories over
several scans. Each plausible assignment extends a hypothesis tree, and
unlikely or redundant branches are pruned as further measurements become
available. MHT can recover from temporary ambiguity, occlusion, and crossing
trajectories because it delays irreversible association decisions. Its main
limitations are memory use, implementation complexity, and computational
growth with the number of targets and plausible assignments
\cite{reid1979algorithm,blackman1999design,barshalom2011tracking}.

% --------------------------------------------
\subsection{Track Initiation, Confirmation, and Deletion}
% --------------------------------------------

Association must be accompanied by track-management logic. An unassociated
detection may indicate a newly appearing object, but it may also be caused
by clutter, multipath, an image artefact, or a transient detector error.
For this reason, new detections usually initialise \emph{tentative} tracks
rather than immediately confirmed objects.

A tentative track may be confirmed after receiving a sufficient number of
consistent updates, for example through an $m$-of-$n$ rule, under which at
least $m$ detections must be obtained during the most recent $n$ scans.
Probabilistic alternatives maintain an explicit track-existence
probability and confirm the track once this probability exceeds a
threshold. Tracks are deleted when their existence probability becomes too
small, when they remain unobserved for a prescribed interval, or when their
predicted state leaves the region of interest
\cite{blackman1999design,barshalom2011tracking}.

These rules determine the responsiveness and stability of the perception
system. Aggressive confirmation detects new road users quickly but may
generate false tracks. Conservative confirmation suppresses false alarms
but may delay the recognition of a newly appearing pedestrian or vehicle.
Likewise, deleting tracks too quickly causes fragmentation during short
occlusions, whereas retaining them too long creates stale or ghost tracks.

% --------------------------------------------
\subsection{Fusion of Independent Tracks}
% --------------------------------------------

Once two tracks have been associated, their state estimates may be fused.
Suppose that
\begin{equation}
\widehat{\mathbf{x}}_{i}=\mathbf{x}+\mathbf{e}_{i},\qquad i\in\{1,2\},
\label{eq:local-track-error}
\end{equation}
where
$\mathbb{E}[\mathbf{e}_{i}]=\mathbf{0}$ and $\operatorname{cov}(\mathbf{e}_{i})=\mathbf{P}_{i}$. If the track errors are independent, their information matrices are
additive, i.e.,  $\mathbf{P}_{\mathrm{f}}^{-1}=\mathbf{P}_{1}^{-1}+\mathbf{P}_{2}^{-1}$. The corresponding fused estimate is
\begin{equation}
\widehat{\mathbf{x}}_{\mathrm{f}}=\mathbf{P}_{\mathrm{f}}\left(\mathbf{P}_{1}^{-1}\widehat{\mathbf{x}}_{1}+\mathbf{P}_{2}^{-1}\widehat{\mathbf{x}}_{2}\right).
\label{eq:independent-track-fusion-state}
\end{equation}
This rule is the track-level counterpart of weighted least squares and information-form Kalman fusion. The weighting is directional rather than scalar. A radar track may dominate the longitudinal-velocity components, while a LiDAR track may dominate lateral position and object extent
\cite{barshalom1995multitarget,barshalom2011tracking}. The independence assumption is nevertheless restrictive. It may be
reasonable when the tracks are produced from completely separate sensors
and have never exchanged information. In practical perception systems,
however, the tracks may share an ego-pose estimate, calibration parameters,
map priors, preprocessing stages, or previously fused observations.
Treating such tracks as independent double-counts their common information
and produces an unrealistically small covariance.

% --------------------------------------------
\subsection{Fusion with Known Cross-Correlation}
% --------------------------------------------

Let the cross-covariance between the local estimation errors be $\operatorname{cov}(\mathbf{e}_{1},\mathbf{e}_{2})=\mathbf{P}_{12}$. When $\mathbf{P}_{12}$ is known, the minimum-variance linear unbiased
fusion rule can be written as
\begin{equation}
\widehat{\mathbf{x}}_{\mathrm{f}}=\widehat{\mathbf{x}}_{1}+\mathbf{W}\left(\widehat{\mathbf{x}}_{2}-\widehat{\mathbf{x}}_{1}\right),
\label{eq:correlated-track-fusion-state}
\end{equation}
where $\mathbf{W}=\left(\mathbf{P}_{1}-\mathbf{P}_{12}\right)\left(\mathbf{P}_{1}+\mathbf{P}_{2}-\mathbf{P}_{12}-\mathbf{P}_{12}^{\mathsf T}\right)^{-1}$. The fused covariance is
\begin{equation}
\mathbf{P}_{\mathrm{f}}=\mathbf{P}_{1}-\mathbf{W}\left(\mathbf{P}_{1}-\mathbf{P}_{12}^{\mathsf T}\right).
\label{eq:correlated-track-fusion-cov}
\end{equation}
When $\mathbf{P}_{12}=\mathbf{0}$, the rule reduces to independent
information fusion. Propagating the cross-covariance gives statistically
efficient estimates, but requires the fusion architecture to preserve the
complete information-sharing history. This bookkeeping becomes difficult
in decentralised, asynchronous, or dynamically changing networks
\cite{barshalom1995multitarget,saha1998track,
barshalom2011tracking}.

% --------------------------------------------
\subsection{Unknown Correlations and Covariance Intersection}
% --------------------------------------------

When the cross-correlation is unknown, covariance intersection (CI)
provides a conservative alternative. Given two consistent estimates
$\widehat{\mathbf{x}}_{1}$ and $\widehat{\mathbf{x}}_{2}$ with covariances
$\mathbf{P}_{1}$ and $\mathbf{P}_{2}$, CI defines
\begin{equation}
\mathbf{P}_{\mathrm{CI}}^{-1}=\omega\mathbf{P}_{1}^{-1}+(1-\omega)\mathbf{P}_{2}^{-1},
\label{eq:ci-cov}
\end{equation}
and
\begin{equation}
\widehat{\mathbf{x}}_{\mathrm{CI}}=\mathbf{P}_{\mathrm{CI}}\left(\omega\mathbf{P}_{1}^{-1}\widehat{\mathbf{x}}_{1}+(1-\omega)\mathbf{P}_{2}^{-1}\widehat{\mathbf{x}}_{2}\right),
\label{eq:ci-est}
\end{equation}
where $\omega\in[0,1]$. The weight is commonly selected by solving
\begin{equation}
\omega^{\star}=\arg\min_{\omega\in[0,1]}J\!\left(\mathbf{P}_{\mathrm{CI}}(\omega)\right),
\label{eq:ci-weight}
\end{equation}
where $J(\mathbf{P})$ may be
$\operatorname{tr}(\mathbf{P})$,
$\det(\mathbf{P})$, or
$\log\det(\mathbf{P})$
\cite{uhlmann1995dynamic,julier1997nondivergent,julier2001general}. CI guarantees consistency without requiring the cross-covariance. It is therefore suitable for merging tracks from independently developed
perception modules, neighbouring vehicles, or roadside infrastructure when their information histories are unavailable. Its price is conservatism: CI avoids claiming that all information contained in the two estimates is independent. When the true cross covariance is known, the correlated fusion rule generally gives a tighter estimate \cite{barshalom2011tracking,julier2001general}. CI should not be interpreted as simple arithmetic averaging. The scalar
$\omega$ weights the complete information matrices. Consequently, the
fusion remains direction dependent: one track may be more informative
about position, while another is more informative about velocity.



% ==============================================================
%
\section{Evidence Theory for Conflicting Sensor Information}
\label{sec:evidence}
%
% ==============================================================

Track-level fusion is well suited to combining numerical state estimates expressed in a common space and characterised by statistical uncertainty. In many systems, however, information may also be qualitative, incomplete, or conflicting, such as class hypotheses, existence assessments, occupancy states, or risk indicators. Evidence theory complements covariance-based fusion by providing a framework for combining such heterogeneous information while explicitly representing ambiguity, ignorance, and disagreement. Autonomous driving provides a representative case study in which these challenges arise across multiple sensing modalities.

Evidence theory, also known as Dempster--Shafer theory, provides a mathematical framework for representing partial belief, ignorance, and conflict between information sources \cite{shafer_mathematical_1976,dempster_upperlower_2008,smets_transferable_1990}. It is particularly useful at the object, track, and decision levels of autonomous-driving fusion, where different modules may produce heterogeneous outputs such as class labels, existence probabilities, occupancy states, or risk scores. For example, a camera may classify a distant object as a pedestrian, LiDAR may support only the weaker hypothesis that an occupied vertical structure is present, and radar may provide a velocity measurement without a reliable class label. Evidence theory allows the fusion module to retain this ambiguity instead of forcing an early single-class decision.

Let the frame of discernment be a finite set of mutually exclusive hypotheses,
\begin{equation}
\Omega = \{\omega_1, \omega_2, \ldots, \omega_n\},
\end{equation}
where the hypotheses may correspond to object classes, occupancy states, sensor-health states, or driving-risk categories. Evidence theory operates on the power set \(\Gamma(\Omega)\) (also commonly written $2^{\Omega}$), namely the set of all subsets of \(\Omega\):
\begin{equation}
\Gamma(\Omega) = \{\gamma_1, \gamma_2, \ldots, \gamma_{2^{|\Omega|}}\}.
\end{equation}
A basic belief assignment (BBA) is a function
\begin{equation}
m: \Gamma(\Omega) \rightarrow [0,1]
\end{equation}
that assigns belief mass to individual hypotheses or to sets of hypotheses, subject in the classical closed-world setting to
\begin{equation}
\sum_{\gamma_a \subseteq \Omega} m(\gamma_a)=1, \qquad m(\emptyset)=0.
\label{eq:bba}
\end{equation}
The subsets \(\gamma_a\) for which \(m(\gamma_a)>0\) are called focal elements. Unlike a probability distribution, a BBA can allocate mass to a non-singleton set, such as \(\{\text{pedestrian},\text{cyclist}\}\), when the available sensor information supports a restricted group of hypotheses but is insufficient to discriminate between them.

Two quantities derived from the BBA are belief and plausibility. For a proposition \(\gamma_a \subseteq \Omega\), belief is the total support committed to \(\gamma_a\),
\begin{equation}
\mathrm{Bel}(\gamma_a)=
\sum_{\gamma_b\subseteq \gamma_a}m(\gamma_b),
\end{equation}
whereas plausibility is the total support that does not contradict it,
\begin{equation}
\mathrm{Pl}(\gamma_a)=
\sum_{\gamma_b\cap \gamma_a\neq \emptyset}m(\gamma_b).
\end{equation}
The interval \([\mathrm{Bel}(\gamma_a),\mathrm{Pl}(\gamma_a)]\) therefore separates confirmed support from remaining possible support. In autonomous driving, this is useful when a perception system needs to report not only that an object is likely to exist, but also how much uncertainty remains because of occlusion, sensor sparsity, adverse weather, or conflicting modality outputs.

When a single operational decision is required, a BBA can also be converted into an ordinary probability distribution using the {\em pignistic} transformation. For a singleton hypothesis \(\omega\in\Omega\), this transformation distributes the mass assigned to each non-singleton focal set equally among the singleton hypotheses contained in that set:
\begin{equation}
\mathrm{BetP}(\omega)=
\sum_{\substack{\gamma_a\subseteq\Omega\\ \omega\in\gamma_a}}
\frac{m(\gamma_a)}{|\gamma_a|}.
\label{eq:pignistic-transform}
\end{equation}
Equation~\eqref{eq:pignistic-transform} holds under the closed-world assumption $m(\emptyset)=0$ adopted here; in the general open-world case, the right-hand side is additionally rescaled by $1/(1-m(\emptyset))$.
The pignistic distribution is therefore useful when an evidence-theoretic fusion module must pass a probability-like score to a downstream classifier, planner, or decision rule, while belief and plausibility still preserve the distinction between committed support and residual uncertainty.

When independent information sources provide BBAs on the same frame of discernment, combination rules can be used to merge them. Dempster's rule emphasises agreement between sources and normalises away the conflicting mass. For \(\gamma_a\neq\emptyset\),
\begin{equation}
(m_i\oplus_D m_j)(\gamma_a)=
\frac{\displaystyle\sum_{\gamma_b \cap \gamma_c=\gamma_a} m_i(\gamma_b)m_j(\gamma_c)}
{1-K},
\label{eq:dempster}
\end{equation}
where
\begin{equation}
K=\sum_{\gamma_b \cap \gamma_c=\emptyset}m_i(\gamma_b)m_j(\gamma_c)
\end{equation}
is the conflict between the two sources. The rule is attractive when the frame of discernment is trusted and the sources are expected to be reliable, but it may produce overconfident results under severe conflict.

The Transferable Belief Model proposed by Smets avoids this normalisation and permits non-zero mass on the empty set \cite{smets_transferable_1990}:
\begin{equation}
(m_i\oplus_S m_j)(\gamma_a)=
\sum_{\gamma_b \cap \gamma_c=\gamma_a}m_i(\gamma_b)m_j(\gamma_c).
\label{eq:smets}
\end{equation}
In this case \((m_i\oplus_S m_j)(\emptyset)=K\), so the empty-set mass explicitly represents contradiction. This interpretation is valuable for safety-critical fusion: a large conflict can indicate sensor malfunction, calibration drift, an unmodelled object class, a stale map prior, or a rapidly changing situation. Recent autonomous-vehicle work has used Dempster--Shafer fusion and evidence-distance weighting for real-time cooperative perception and collision warning, illustrating the relevance of evidence-theoretic reasoning beyond purely theoretical uncertainty modelling \cite{cho_realtime_2020}.

In practical systems, BBAs are obtained from calibrated classifier scores, detector confidences, occupancy-grid likelihoods, measurement residuals, or sensor-health indicators. A common implementation first discounts each source according to its current reliability \(\alpha_i\in[0,1]\):
\begin{equation}
m_i^{\alpha_i}(\gamma_a)=\alpha_i m_i(\gamma_a),
\qquad \gamma_a \subset \Omega,
\end{equation}
\begin{equation}
m_i^{\alpha_i}(\Omega)=1-\alpha_i+\alpha_i m_i(\Omega).
\end{equation}
The discounted BBAs are then combined. In autonomous driving, \(\alpha_i\) may depend on visibility, LiDAR point density, radar signal-to-noise ratio, calibration residuals, timestamp delay, map freshness, or historical sensor consistency. This connects evidence theory naturally with reliability-aware fusion and multi-criteria decision-making methods, which can be used to estimate the discount factors from several competing quality indicators.

\begin{figure}[t]
    \centering
    \begin{tikzpicture}[
        font=\scriptsize,
        box/.style={draw=blue!70!black, rounded corners=2pt, fill=blue!5, align=center, text width=2.25cm, minimum height=1.05cm},
        fuse/.style={draw=blue!85!black, rounded corners=2pt, fill=blue!10, align=center, text width=2.1cm, minimum height=1.0cm},
        decision/.style={draw=blue!95!black, rounded corners=2pt, fill=blue!15, align=center, text width=2.25cm, minimum height=1.05cm},
        >=Latex
    ]
        \node[box] (cam) at (0,0) {Camera\\glare + occlusion\\\(m_C(F)=0.55\)};
        \node[box] (rl) at (0,-1.35) {Radar--LiDAR\\persistent return\\\(m_{RL}(O)=0.45\)};
        \node[fuse] (comb) at (3.05,-0.675) {Evidence\\combination\\conflict \(K=0.44\)};
        \node[decision] (plan) at (6.0,-0.675) {Planner signal\\\(\mathrm{Pl}(O)=0.80\)\\slow + rescan};
        \draw[->, thick, blue!75!black] (cam.east) -- (comb.west);
        \draw[->, thick, blue!75!black] (rl.east) -- (comb.west);
        \draw[->, thick, blue!75!black] (comb.east) -- (plan.west);
    \end{tikzpicture}
    \caption{Evidence-theoretic fusion for a rain-night autonomous-driving conflict. The conflict mass is not merely a nuisance term: it is a safety signal that the perception stack should request additional sensing or choose a conservative manoeuvre.}
    \label{fig:evidence-theory-autonomous-driving-example}
\end{figure}
\begin{example}[Rain-night occlusion conflict.]
Consider an autonomous vehicle approaching a poorly illuminated pedestrian crossing during heavy rain. The camera stream is affected by glare and partial occlusion, while the radar--LiDAR track-level module returns a weak but persistent object hypothesis. Let
\(\Omega=\{F,O,U\}\), where \(F\) denotes a free lane, \(O\) denotes an occupied lane caused by a vehicle or obstacle, and \(U\) denotes an unknown vulnerable road user or unmodelled hazard. Suppose the camera-derived BBA \(m_C\) and the radar--LiDAR BBA \(m_{RL}\) are
\begin{equation}
\begin{array}{c|cccc}
 & F & O & \{O,U\} & \Omega \\
\hline
m_C & 0.55 & 0 & 0.20 & 0.25 \\
m_{RL} & 0 & 0.45 & 0.35 & 0.20 .
\end{array}
\label{eq:evidence-example-bbas}
\end{equation}
The camera is therefore leaning toward ``free lane'', but it still admits ambiguity; the radar--LiDAR module instead supports a non-free lane without being able to decide whether the cause is an ordinary obstacle or a vulnerable road user.

Using the unnormalised Smets rule, the empty-set mass is
\(K=m_S(\emptyset)=0.55(0.45+0.35)=0.44\). If Dempster's normalised rule is used, the fused BBA becomes
\begin{equation}
\begin{aligned}
m_D(F)&=0.1964, &
m_D(O)&=0.3616,\\
m_D(\{O,U\})&=0.3527, &
m_D(\Omega)&=0.0893 .
\end{aligned}
\end{equation}
Consequently,
\begin{equation}
\begin{aligned}
\mathrm{Bel}(O)&=0.3616,\\
\mathrm{Pl}(O)&=0.3616+0.3527+0.0893=0.8036 .
\end{aligned}
\end{equation}
Applying the pignistic transformation~\eqref{eq:pignistic-transform} to the fused BBA gives approximately \(\mathrm{BetP}(F)=0.2262\), \(\mathrm{BetP}(O)=0.5677\), and \(\mathrm{BetP}(U)=0.2061\). Thus, although the most likely singleton is an obstacle, the system should retain a meaningful unknown-road-user component. In an autonomous-driving stack this is exactly the kind of case where evidence theory is useful: the fused output can feed the planner as a cautious risk estimate, while the conflict \(K=0.44\) separately triggers a rescan, reduced speed, or a request for cooperative perception. 
\end{example}
\smallskip
The main limitation of evidence theory is computational complexity: the power set grows exponentially with the number of hypotheses. Practical autonomous-driving applications therefore keep the frame of discernment small, group rare classes into an \emph{unknown} or \emph{other} category, approximate the set of focal elements, or apply evidence-theoretic reasoning only at selected high-value decision points such as conflict detection, object-existence fusion, free-space validation, cooperative perception, and sensor-health monitoring.

A related extension is Dezert--Smarandache Theory (DSmT), which relaxes the strict mutual-exclusivity assumption of classical Dempster--Shafer theory \cite{smarandache_intro_dsmt_2009}. Instead of assigning masses only on the power set of exclusive hypotheses, DSmT can operate on a richer fusion space that also includes intersections between hypotheses. This is useful when the frame cannot be cleanly refined into mutually exclusive classes, for example when an autonomous-driving system must reason about overlapping or ambiguous propositions such as a partially observed object that may simultaneously support \emph{pedestrian-like} and \emph{cyclist-like} evidence, or when different sensors produce paradoxical but informative observations. The trade-off is that exact DSmT can be computationally demanding. Therefore, its most practical role in real-time autonomous driving is likely to be through restricted models, reduced focal sets, or approximate DSmT combination rules designed to reduce computational cost while preserving the ability to reason under strong conflict \cite{guo_evidenceclustering_2015}.
\section{Optimisation-Based Fusion and Factor Graphs}
\label{sec:factor-graphs}

Optimisation-based fusion estimates the most likely state or trajectory by minimising an objective function that combines multiple measurement residuals. In autonomous driving, optimisation is central to visual--inertial odometry, LiDAR--inertial odometry, GNSS--IMU fusion, map matching, bundle adjustment, and pose graph optimisation.

A typical formulation defines a state vector containing poses, velocities, sensor biases, map landmarks, or object states. Each sensor measurement contributes a residual term that penalises the discrepancy between the predicted measurement and the observed measurement. The final estimate is obtained by minimising a weighted sum of residuals, where the weights are related to measurement uncertainty. This formulation naturally supports heterogeneous sensors because each measurement type can be represented as a factor connected to the relevant state variables.

Factor graphs provide a structured representation of this problem. Nodes represent unknown states, while factors represent constraints from IMU integration, camera reprojection, LiDAR scan matching, GNSS measurements, wheel odometry, loop closures, or map priors. Compared with purely filtering-based methods, smoothing and factor graph approaches can use measurements over a time window or an entire trajectory, which can improve consistency and accuracy. Their main limitation is computational cost, particularly for real-time deployment on embedded platforms. Standard references and tooling have made this formulation the de facto back-end of modern localisation stacks: iSAM2 introduced incremental smoothing via the Bayes tree, and GTSAM and related libraries provide reusable factor abstractions \cite{kaess2012_isam2,dellaert2017_factor_graphs}. % [July 2026 pass] This subsection previously cited nothing; added the two canonical factor-graph references.

\section{Fusion based on Multi-Criteria Decision-Making}
\label{sec:mcdm}

Multi-criteria decision-making (MCDM) provides a family of methods for ranking alternatives when several criteria must be considered simultaneously. In sensor fusion, the alternatives may be sensors, sensor suites, object hypotheses, fusion algorithms, tracks, or possible system actions. The criteria may include detection accuracy, localisation error, latency, energy consumption, robustness to weather, calibration quality, field of view, uncertainty, interpretability, cost, and safety relevance. MCDM is therefore not a replacement for probabilistic filtering or deep multimodal fusion; rather, it is a decision-support layer that can select, weight, or audit fusion choices when several imperfect sources must be balanced.

The discussion below focuses mainly on the Analytic Hierarchy Process (AHP), because pairwise comparisons provide a natural way to elicit expert preferences and sensor-reliability weights. However, AHP is only one member of a broader MCDM family. Outranking methods such as ELECTRE III can compare alternatives through concordance, discordance, and veto thresholds when strict scalar scoring is too reductive \cite{roy_electreiii_1978}. DEMATEL can instead be used to model causal influence between criteria, for example when weather affects visibility, visibility affects detector confidence, and confidence affects the trust assigned to a perception module \cite{gabus_fontela_dematel_1972}. Such methods can be useful complements when fusion criteria are interdependent or when the goal is not simply to compute a single weighted sum.

This perspective is useful in autonomous driving because the most reliable sensing modality depends on the context. Cameras are strong for semantics but weak under glare or darkness. LiDAR is strong for geometry but may lose range and density in fog, rain, or snow. Radar is robust to poor visibility and provides direct velocity, but its angular resolution is limited. A fusion system can use MCDM to compute context-dependent reliability weights before combining sensor outputs, or to choose between alternative perception pipelines under resource constraints.

Let \(\mathcal{A}=\{a_1,\ldots,a_r\}\) be a set of alternatives and \(\mathcal{C}=\{c_1,\ldots,c_p\}\) a set of criteria. After normalisation, the performance of alternative \(a_i\) with respect to criterion \(c_k\) can be represented as \(q_{ik}\in[0,1]\). Given criterion weights \(w_k\geq 0\), \(\sum_{k=1}^{p} w_k=1\), a simple additive MCDM score is
\begin{equation}
S(a_i)=\sum_{k=1}^{p} w_k q_{ik}.
\label{eq:mcdm-score}
\end{equation}
For sensor fusion, \(S(a_i)\) may represent the suitability of sensor \(i\), the reliability of a detector, or the priority of a candidate response. The resulting scores can be used directly for selection or transformed into fusion weights.

AHP is a widely used MCDM method in which criterion weights are elicited through pairwise comparisons \cite{saaty_analytic_1987}. In the fusion setting, an expert may compare criteria such as detection confidence, latency, adverse-weather robustness, and calibration quality. A pairwise comparison \(A_{ij}\) expresses how strongly criterion \(i\) is preferred to criterion \(j\), typically using Saaty's scale, with reciprocal values satisfying \(A_{ji}=1/A_{ij}\). For complete and consistent comparisons, the priority vector can be estimated from the principal eigenvector of the comparison matrix. Sparse and distributed AHP extensions are especially relevant when only selected pairwise comparisons are available, as is common when experts compare many sensors, criteria, or failure modes \cite{oliva_sparse_distributed_2017,bozoki_fulop_ronyai_2010,bozoki_optimality_2019}.

Consider \(g\) decision-makers or design sources. Each decision-maker \(u\) provides a possibly incomplete set of comparisons represented by a graph \(\mathcal{G}^{(u)}=(\mathcal{V},\mathcal{E}^{(u)})\). The node set \(\mathcal{V}=\{1,\ldots,r\}\) represents the alternatives or criteria to be ranked, while an edge \((i,j)\in\mathcal{E}^{(u)}\) means that decision-maker \(u\) has provided the comparison \(A_{ij}^{(u)}\). Because pairwise comparison matrices are reciprocal, the reverse comparison is implicit: \(A_{ji}^{(u)}=1/A_{ij}^{(u)}\). Let \(y_i=\ln x_i\), where \(x_i>0\) is the unknown priority of alternative or criterion \(i\). The group logarithmic least-squares (LLS) problem can be written as
\begin{equation}
\bm{y}^{*}=
\arg\min_{\bm{y}\in\mathbb{R}^{r}}
\frac{1}{2}
\sum_{u=1}^{g}
\sum_{(i,j)\in\mathcal{E}^{(u)}}
\left(\ln A_{ij}^{(u)}-y_i+y_j\right)^2 .
\label{eq:lls-ahp}
\end{equation}
The inner sum therefore runs only over the observed comparisons, and each residual measures the mismatch between the expert's logarithmic judgement \(\ln A_{ij}^{(u)}\) and the logarithmic ratio \(y_i-y_j=\ln(x_i/x_j)\) implied by the estimated priorities.

The problem in \eqref{eq:lls-ahp} has an intrinsic scaling ambiguity. Adding the same constant to all components of \(\bm{y}\) does not change any difference \(y_i-y_j\), and multiplying all positive priorities by the same constant does not change any ratio \(x_i/x_j\). Thus, when the collective comparison graph is connected, the solution is unique only up to an additive constant in \(\bm{y}\), equivalently up to a multiplicative scale factor in \(\bm{x}\). One can fix this ambiguity by imposing a gauge such as \(y_1=0\) or \(\sum_i y_i=0\), and then normalise the positive priorities as
\begin{equation}
x_i^{*}=\frac{\exp(y_i^{*})}{\sum_{\ell=1}^{r}\exp(y_\ell^{*})}.
\end{equation}
This formulation is useful because it can handle missing pairwise comparisons and multiple experts. If the objective aggregates all decision-makers into one problem, it produces a single collective priority vector. If the expert population is heterogeneous, the same pairwise-comparison structure can also be used to identify clusters of decision-makers with internally coherent opinions and mutually different rankings \cite{oliva_opinion_2019}.

\smallskip
\begin{example}[Incomplete AHP for a rain-night sensor policy.]
As a small autonomous-driving example, suppose a fallback supervisor must rank four information streams for the same adverse-weather scene: camera semantics \(C\), LiDAR geometry \(L\), radar range--rate \(R\), and a map/IMU lane prior \(M\). Expert comparisons are available only for a sparse connected graph, not for the full \(4\times4\) matrix:
\[
\begin{aligned}
A_{L,C}&=2, & A_{R,C}&=5,\\
A_{M,C}&=3, & A_{R,M}&=1.5 .
\end{aligned}
\]
The missing comparisons \(A_{R,L}\) and \(A_{M,L}\) are not filled by hand. Instead, the LLS problem in \eqref{eq:lls-ahp} estimates the log-priorities from the observed edges only. With the gauge \(\sum_i y_i=0\), the solution gives the normalised priority vector
\begin{equation}
\begin{aligned}
\bm{x}^{*}
&=\begin{bmatrix}x_C & x_L & x_R & x_M\end{bmatrix}\\
&=\begin{bmatrix}0.091 & 0.183 & 0.441 & 0.284\end{bmatrix}.
\end{aligned}
\label{eq:incomplete-ahp-example-weights}
\end{equation}
The ranking is therefore \(R \succ M \succ L \succ C\): radar receives the largest weight because range--rate robustness is highly valuable in rain and darkness, while the map/IMU prior remains important for lane-level consistency even though it does not directly observe the obstacle. The same solution also predicts the two omitted ratios,
\begin{equation}
\widehat{A}_{R,L}\approx2.41,\qquad
\widehat{A}_{M,L}\approx1.55,
\end{equation}
which can be shown to the expert as a consistency check before deployment.

\begin{figure}[t]
    \centering
    \begin{tikzpicture}[
        font=\scriptsize,
        node/.style={draw=blue!75!black, rounded corners=2pt, fill=blue!5, align=center, text width=1.55cm, minimum height=0.8cm},
        obs/.style={->, thick, blue!75!black},
        miss/.style={->, thick, dashed, red!45!red},
        >=Latex
    ]
        \node[node] (C) at (0,0) {Camera\\\(0.091\)};
        \node[node] (L) at (2.35,1.35) {LiDAR\\\(0.183\)};
        \node[node] (R) at (4.7,0) {Radar\\\(0.441\)};
        \node[node] (M) at (2.35,-1.35) {Map/IMU\\\(0.284\)};

        \draw[obs] (L) -- node[above left, blue] {\(2\)} (C);
        \draw[obs] (R) -- node[above left, blue] {\(5\)} (C);
        \draw[obs] (M) -- node[below left, blue] {\(3\)} (C);
        \draw[obs] (R) -- node[below right, blue] {\(1.5\)} (M);

        \draw[miss] (R) -- node[above right, red] {\(2.41\)} (L);
        \draw[miss] (M) -- node[below right, red] {\(1.55\)} (L);
    \end{tikzpicture}
    \caption{Incomplete AHP graph for context-dependent sensor-stream reliability in autonomous driving. Solid arrows are elicited pairwise comparisons; dashed arrows are ratios inferred by the logarithmic least-squares completion.}
    \label{fig:incomplete-ahp-autonomous-driving-example}
\end{figure}

These weights can be used directly as context-dependent sensor-selection weights, or they can be converted into evidence-discount factors before applying \eqref{eq:dempster} or \eqref{eq:smets}. In the rain-night example, for instance, the low camera weight does not delete the camera information; it simply prevents a visually overconfident ``free-lane'' claim from dominating radar and LiDAR evidence. 
\end{example}
\smallskip

A further issue is robustness. Pairwise comparisons are subjective and can be affected by uncertainty, bias, or simple data-entry errors. In sensor fusion, this matters because small changes in expert judgements may alter the ranking of sensing modalities or fallback strategies. Robustness to rank reversal can be analysed by searching for the smallest perturbations of the comparison values that would change the ordering of at least two alternatives \cite{faramondi_robustness_2023}. Related work studies order preservation, inconsistency, incomplete comparison matrices, and the effect of different weighting procedures, providing useful checks on whether an AHP-derived ranking should be trusted \cite{kulakowski_orderpreservation_2015,agoston_csato_inconsistency_2022,csato_agoston_bozoki_2024}. Recent robust group-decision work also considers experts whose judgements are outlying or manipulative and proposes heuristic aggregation rules that reduce their influence \cite{kulakowski_resilient_2024}.

MCDM can also be integrated with evidence-theoretic fusion. For example, each sensor can be assessed according to several reliability criteria: current measurement quality, historical false-alarm rate, conflict with other sensors, environmental suitability, latency, and calibration residual. The resulting MCDM score can define a discount factor \(\alpha_i\) for the BBA of sensor \(i\), after which the discounted evidence is combined using Dempster's rule, Smets' rule, or a weighted evidence rule. This type of approach has been studied in multi-sensor fusion based on belief functions, where AHP is used to weight uncertain and conflicting sensor readings before combination \cite{frikha_ahp_2015}.

In autonomous driving, MCDM-based fusion is most appropriate when the decision problem is explicit and the criteria are interpretable. Examples include selecting the most reliable modality under degraded sensing, ranking candidate sensor suites during system design, choosing whether to trust an online map prior, selecting a fallback perception pipeline under compute constraints, or prioritising collision-warning responses. Its main advantages are transparency and auditability: the system can explain why a sensor was down-weighted or why a conservative action was selected. Its limitations are the subjectivity of criterion weights, sensitivity to normalisation choices, possible rank changes under perturbations, and the need to validate expert-driven rules under rare safety-critical scenarios.

\problemsection
Solutions are collected at the end of the book. Every number used below was obtained with the scripts of the companion repository, \repopath{companion\_\allowbreak repo/\allowbreak exercises/\allowbreak ch03/\allowbreak }, which also contains the problem statements in machine-readable form and the Jupyter notebooks of the computational problems.

\begin{prob}\label{prob:ch03:1}
\textbf{(Least squares versus weighted least squares.)} A radar and a camera both report the longitudinal distance $x$ to a lead vehicle: $y_1=41.6$\,m with $\sigma_1=0.30$\,m and $y_2=43.4$\,m with $\sigma_2=1.50$\,m, the two errors being zero-mean and independent.
\begin{enumerate}
\item[(a)] Write the stacked linear model \eqref{eq:linear-measurement-model}, i.e.\ $\mathbf{G}=\begin{bmatrix}1&1\end{bmatrix}^{\mathsf T}$ and $\mathbf{R}=\diag(\sigma_1^{2},\sigma_2^{2})$, and compute the ordinary least-squares estimate $\widehat{x}_{\mathrm{LS}}$ together with its variance $\operatorname{cov}(\widehat{x}_{\mathrm{LS}})=(\mathbf{G}^{\mathsf T}\mathbf{G})^{-1}\mathbf{G}^{\mathsf T}\mathbf{R}\mathbf{G}(\mathbf{G}^{\mathsf T}\mathbf{G})^{-1}$.
\item[(b)] Compute the information $\mathbf{\Lambda}=\mathbf{G}^{\mathsf T}\mathbf{R}^{-1}\mathbf{G}$, the WLS estimate $\widehat{x}_{\mathrm{WLS}}$ and its variance $\mathbf{\Lambda}^{-1}$.
\item[(c)] Show that, in general, $\widehat{x}_{\mathrm{WLS}}=(\sigma_1^{-2}y_1+\sigma_2^{-2}y_2)/(\sigma_1^{-2}+\sigma_2^{-2})$, i.e.\ the fused estimate is the information-weighted average of the two readings, and that its variance never exceeds the variance of the better sensor alone.
\item[(d)] By which factor does ordinary least squares inflate the variance in this example, and why is the WLS estimate so close to the radar reading?
\end{enumerate}
\end{prob}

\begin{prob}\label{prob:ch03:2}
\textbf{(Regularised least squares as prior-plus-measurement fusion.)} A tracker predicts the range and range rate of the same lead vehicle as $\mathbf{m}=\begin{bmatrix}42.0 & -5.0\end{bmatrix}^{\mathsf T}$ (m, m/s) with prior covariance
$\mathbf{P}=\begin{bmatrix}1.0&0.3\\0.3&0.25\end{bmatrix}$.
The radar of Problem~\ref{prob:ch03:1} then measures the range only, $y=41.6$\,m with $\mathbf{R}=0.09$\,m$^{2}$, so that $\mathbf{G}=\begin{bmatrix}1&0\end{bmatrix}$.
\begin{enumerate}
\item[(a)] Compute $\mathbf{P}^{-1}$, the posterior information matrix $\mathbf{\Lambda}_{\mathrm{post}}=\mathbf{G}^{\mathsf T}\mathbf{R}^{-1}\mathbf{G}+\mathbf{P}^{-1}$ and the posterior covariance $\mathbf{P}^{+}$.
\item[(b)] Compute $\widehat{\mathbf{x}}_{\mathrm{ReLS}}=\widehat{\mathbf{x}}_{\mathrm{MAP}}$.
\item[(c)] Verify the result through the correction form of Section~\ref{subsubsec:sequential_least_squares}, $\mathbf{K}=\mathbf{P}\mathbf{G}^{\mathsf T}(\mathbf{G}\mathbf{P}\mathbf{G}^{\mathsf T}+\mathbf{R})^{-1}$, $\widehat{\mathbf{x}}=\mathbf{m}+\mathbf{K}(y-\mathbf{G}\mathbf{m})$, $\mathbf{P}^{+}=(\mathbf{I}-\mathbf{K}\mathbf{G})\mathbf{P}$.
\item[(d)] The range rate is not measured, yet it is corrected. Explain why, and state what happens when the off-diagonal entry of $\mathbf{P}$ is set to zero.
\item[(e)] Discuss the limits $\mathbf{P}^{-1}\to\mathbf{0}$ (no prior information) and $\mathbf{R}^{-1}\to\mathbf{0}$ (useless measurement).
\end{enumerate}
\end{prob}

\begin{prob}\label{prob:ch03:3}
\textbf{(Sequential least squares is batch least squares.)} Two measurements $\mathbf{y}_1=\mathbf{G}_1\mathbf{x}+\mathbf{r}_1$ and $\mathbf{y}_2=\mathbf{G}_2\mathbf{x}+\mathbf{r}_2$, with independent zero-mean noises of covariances $\mathbf{R}_1$ and $\mathbf{R}_2$, are available together with a prior $(\mathbf{m},\mathbf{P}_0)$ on the static unknown $\mathbf{x}$.
\begin{enumerate}
\item[(a)] Apply the information-form update of Section~\ref{subsubsec:sequential_least_squares} twice and show that the resulting estimate and covariance coincide with the batch regularised estimate obtained from the stacked matrix $\mathbf{G}=\begin{bmatrix}\mathbf{G}_1^{\mathsf T}&\mathbf{G}_2^{\mathsf T}\end{bmatrix}^{\mathsf T}$ and the block-diagonal covariance $\mathbf{R}=\diag(\mathbf{R}_1,\mathbf{R}_2)$.
\item[(b)] Deduce that the result does not depend on the order in which the two measurements are processed.
\item[(c)] Take the radar of Problem~\ref{prob:ch03:1} as the starting point ($\widehat{x}_1=41.6$, $P_1=0.09$) and process the camera reading with the gain form; check numerically that the WLS values of Problem~\ref{prob:ch03:1} are recovered.
\item[(d)] Identify the step of the derivation that fails when $\mathbf{r}_1$ and $\mathbf{r}_2$ are correlated, and relate it to the discussion of double counting in Section~\ref{subsec:track-fusion}.
\end{enumerate}
\end{prob}

\begin{prob}\label{prob:ch03:4}
\textbf{(One Kalman-filter step by hand.)} A lead vehicle is tracked along one dimension with the state $\mathbf{x}=\begin{bmatrix}p&v\end{bmatrix}^{\mathsf T}$ (relative position and velocity), the constant-velocity model $\mathbf{F}=\begin{bmatrix}1&T\\0&1\end{bmatrix}$ with $T=0.5$\,s and the discretised white-noise-acceleration covariance $\mathbf{Q}=\sigma_a^{2}\begin{bmatrix}T^{4}/4&T^{3}/2\\T^{3}/2&T^{2}\end{bmatrix}$ with $\sigma_a=0.8$\,m/s$^{2}$, as in Example~\ref{ex:ekf-radar-cutin}. At time $k-1$ the posterior is $\widehat{\mathbf{x}}_{k-1|k-1}=\begin{bmatrix}50.0&-4.0\end{bmatrix}^{\mathsf T}$, $\mathbf{P}_{k-1|k-1}=\diag(0.25,0.16)$. At time $k$ a range-only measurement $y_k=47.4$\,m arrives, with $\mathbf{H}=\begin{bmatrix}1&0\end{bmatrix}$ and $\mathbf{R}=0.09$\,m$^{2}$.
\begin{enumerate}
\item[(a)] Compute $\mathbf{Q}$.
\item[(b)] Compute the prediction \eqref{eq:kf_state_prediction}--\eqref{eq:kf_cov_prediction}.
\item[(c)] Compute the innovation \eqref{eq:kf_innovation}, its covariance \eqref{eq:kf_innovation_covariance} and the gain \eqref{eq:kf_gain}.
\item[(d)] Compute the posterior \eqref{eq:kf_state_update}--\eqref{eq:kf_cov_update} and verify that the Joseph form of Section~\ref{subsubsec:sequential_least_squares} gives the same covariance.
\item[(e)] Compute the NIS \eqref{eq:nis} and compare it with the two-sided $95\,\%$ interval of the $\chi^{2}_{1}$ distribution, $[0.00098,\,5.02]$.
\item[(f)] Only the position is measured, yet the velocity is corrected and the position standard deviation drops from $0.55$\,m (predicted) to $0.26$\,m. Explain both facts in terms of the entries of $\mathbf{P}_{k|k-1}$.
\end{enumerate}
\end{prob}

\begin{prob}\label{prob:ch03:5}
\textbf{(EKF Jacobian of a range--bearing sensor and azimuth wrapping.)} A sensor observes $\mathbf{g}(\mathbf{x})=\begin{bmatrix}\sqrt{p_x^{2}+p_y^{2}} & \operatorname{atan2}(p_y,p_x)\end{bmatrix}^{\mathsf T}$ of the state $\mathbf{x}=\begin{bmatrix}p_x&p_y&v_x&v_y\end{bmatrix}^{\mathsf T}$, with $\sigma_\theta=0.5^{\circ}$.
\begin{enumerate}
\item[(a)] Derive the Jacobian $\mathbf{H}=\partial\mathbf{g}/\partial\mathbf{x}$ and check that it coincides with the first two rows of \eqref{eq:ekf-radar-jacobian}.
\item[(b)] Evaluate it at $\widehat{\mathbf{x}}_{k|k-1}=\begin{bmatrix}40&3&-5&0\end{bmatrix}^{\mathsf T}$.
\item[(c)] The bearing row scales as $1/d$. Explain the consequence for the lateral position information and compute the lateral standard deviation implied by $\sigma_\theta$ at $60$\,m and at $10$\,m.
\item[(d)] The predicted bearing is $3.135$\,rad and the sensor returns $-3.140$\,rad. Compute the raw difference and the innovation wrapped to $(-\pi,\pi]$. Assuming that the innovation covariance is dominated by the sensor noise, $S\approx\sigma_\theta^{2}$, compute the NIS in both cases and state what a $\chi^{2}_{1}$ gate at the $97.5\,\%$ quantile ($5.02$) would do with the detection.
\item[(e)] Explain why the wrapping must be applied to the innovation and cannot be replaced by wrapping the predicted measurement or the state.
\end{enumerate}
\end{prob}

\begin{prob}\label{prob:ch03:6}
\textbf{(Zonotope propagation and a single-strip correction.)} A scalar position and velocity are enclosed at time $k$ by the zonotope $\mathcal{X}_k=p\oplus H\mathbb{B}_{\infty}^{2}$ with
$p=\begin{bmatrix}1\\0\end{bmatrix}$, $H=\begin{bmatrix}1&0.5\\0&1\end{bmatrix}$.
The dynamics are the constant-velocity model $F=\begin{bmatrix}1&1\\0&1\end{bmatrix}$ ($T=1$\,s) driven by an unknown acceleration bounded by $\bar a=0.2$\,m/s$^{2}$, which maps onto the generator $\bar a\begin{bmatrix}T^{2}/2&T\end{bmatrix}^{\mathsf T}$ as in Example~\ref{ex:zono-radar-cutin}.
\begin{enumerate}
\item[(a)] Compute the predicted zonotope $\widetilde{\mathcal{X}}_{k+1}=\tilde p\oplus\tilde H\mathbb{B}_{\infty}^{3}$ and its interval hull.
\item[(b)] The position is then measured as $y=1.4$ with a bounded error $|\Delta y|\le\bar\phi=0.3$. Write the strip $\mathcal{S}$ with normal $c=\begin{bmatrix}1&0\end{bmatrix}^{\mathsf T}$.
\item[(c)] Compute the gain $\lambda^{\star}$ of \eqref{eq:zono-optimal-gain} (the strip gain $\gamma$ of Table~\ref{tab:notation}), the corrected centre $\hat p$, the corrected generator matrix $\hat H$ (four columns) and the resulting interval hull.
\item[(d)] Repeat with the choice $\lambda=c$ used in Example~\ref{ex:zonotope-strip}, and compare the two enclosures through their interval hulls and through $\|\hat H\|_{F}^{2}$. Which quantity does each choice minimise?
\item[(e)] Compute the exact interval hull of $\widetilde{\mathcal{X}}_{k+1}\cap\mathcal{S}$ (two linear programmes over $\xi\in[-1,1]^{3}$ suffice) and comment on the conservatism of both zonotopic enclosures.
\end{enumerate}
\end{prob}

\begin{prob}\label{prob:ch03:7}
\textbf{(Dempster's rule versus Smets' rule.)} Let $\Omega=\{P,V,S\}$ (pedestrian, vehicle, static obstacle). A camera classifier and a LiDAR shape classifier return the BBAs
\begin{equation*}
\begin{array}{c|ccccc}
 & P & V & S & \{P,V\} & \Omega\\
\hline
m_C & 0.5 & 0 & 0 & 0.3 & 0.2\\
m_L & 0 & 0.4 & 0.1 & 0.4 & 0.1
\end{array}
\end{equation*}
\begin{enumerate}
\item[(a)] Compute the conflict mass $K$.
\item[(b)] Compute the fused BBA under Dempster's rule \eqref{eq:dempster}.
\item[(c)] Compute $\mathrm{Bel}$, $\mathrm{Pl}$ and the pignistic probability \eqref{eq:pignistic-transform} of each singleton.
\item[(d)] Compute the fused BBA under Smets' rule \eqref{eq:smets}, including $m(\emptyset)$, and explain in one sentence what the two rules do differently with the same products.
\item[(e)] Two other sources return $m_1(P)=0.99$, $m_1(S)=0.01$ and $m_2(V)=0.99$, $m_2(S)=0.01$. Compute $K$ and the two combinations, and explain why the Dempster result is unsafe for a planner.
\end{enumerate}
\end{prob}

\begin{prob}\label{prob:ch03:8}
\textbf{(Covariance intersection.)} Two tracks $(\widehat{\mathbf{x}}_1,\mathbf{P}_1)$ and $(\widehat{\mathbf{x}}_2,\mathbf{P}_2)$ of the same object are consistent, i.e.\ their true error covariances $\bar{\mathbf{P}}_1$ and $\bar{\mathbf{P}}_2$ satisfy $\bar{\mathbf{P}}_i\preceq\mathbf{P}_i$, but their cross-covariance $\bar{\mathbf{P}}_{12}$ is unknown.
\begin{enumerate}
\item[(a)] Prove that the CI estimate \eqref{eq:ci-cov}--\eqref{eq:ci-est} is consistent for \emph{every} $\omega\in[0,1]$ and every $\bar{\mathbf{P}}_{12}$ for which the joint error covariance is positive semidefinite. \emph{Hint:} write the CI error as $\mathbf{e}=\mathbf{P}_{\mathrm{CI}}\left(\omega\mathbf{P}_1^{-1}\mathbf{e}_1+(1-\omega)\mathbf{P}_2^{-1}\mathbf{e}_2\right)$ and show first that
\begin{equation*}
\begin{bmatrix}\bar{\mathbf{P}}_1&\bar{\mathbf{P}}_{12}\\ \bar{\mathbf{P}}_{12}^{\mathsf T}&\bar{\mathbf{P}}_2\end{bmatrix}\preceq\begin{bmatrix}\bar{\mathbf{P}}_1/\omega&\mathbf{0}\\ \mathbf{0}&\bar{\mathbf{P}}_2/(1-\omega)\end{bmatrix}.
\end{equation*}
\item[(b)] Let $\mathbf{P}_1=\diag(0.25,\,1.0)$ (radar track, precise longitudinally) and $\mathbf{P}_2=\diag(1.0,\,0.16)$ (camera track, precise laterally). Write $\tr\mathbf{P}_{\mathrm{CI}}(\omega)$ in closed form, find the weight $\omega^{\star}$ that minimises it by solving the stationarity condition, and compute $\mathbf{P}_{\mathrm{CI}}(\omega^{\star})$.
\item[(c)] Compare with the covariance of the independent-track fusion rule of Section~\ref{subsec:track-fusion}, $\mathbf{P}_{\mathrm{f}}^{-1}=\mathbf{P}_1^{-1}+\mathbf{P}_2^{-1}$, and explain why the CI result remains direction dependent although $\omega$ is a scalar.
\end{enumerate}
\end{prob}

\begin{prob}\label{prob:ch03:9}
\textbf{(Incomplete AHP by logarithmic least squares.)} A single expert compares four reliability criteria, labelled $1$ to $4$, but only on the pairs
$A_{12}=3$, $A_{23}=2$, $A_{13}=4$, $A_{34}=2$.
\begin{enumerate}
\item[(a)] Draw the comparison graph. Is it connected? Is the cycle $1\!-\!2\!-\!3$ consistent?
\item[(b)] Show that the normal equations of \eqref{eq:lls-ahp} read $\mathbf{L}\bm{y}=\bm{b}$, where $\mathbf{L}$ is the Laplacian of the undirected comparison graph and $b_i=\sum_{j}\ln A_{ij}$ over the neighbours of $i$, with $A_{ji}=1/A_{ij}$; write $\mathbf{L}$ and $\bm{b}$ explicitly.
\item[(c)] Solve with the gauge $\sum_i y_i=0$ and normalise the priorities $x_i^{*}$.
\item[(d)] Compute the fitted ratios $x_i^{*}/x_j^{*}$ and the residuals on the four edges. How is the cycle inconsistency $\ln(A_{12}A_{23}/A_{13})$ distributed, and why is the edge $(3,4)$ fitted exactly?
\item[(e)] Predict the missing comparisons $A_{14}$ and $A_{24}$.
\item[(f)] What changes if the comparison $A_{13}$ is removed?
\end{enumerate}
\end{prob}

\begin{prob}\label{prob:ch03:10}
\textbf{(Computational: process noise and innovation consistency.)} Reproduce the radar cut-in scenario of Example~\ref{ex:ekf-radar-cutin} (ego vehicle at $27$\,m/s, target at $22$\,m/s initially $60$\,m ahead in the adjacent lane, lane change from $3.5$\,m to $0$\,m between $t=5$\,s and $t=8$\,s, radar at $20$\,Hz with $\sigma_r=0.30$\,m, $\sigma_\theta=0.5^{\circ}$, $\sigma_{\dot r}=0.10$\,m/s) and run the constant-velocity EKF on one and the same detection sequence for $\sigma_a\in\{0.1,0.2,0.4,0.8,1.6,3.2,6.4\}$\,m/s$^{2}$. For each value report the position RMSE, the mean NIS \eqref{eq:nis} and the fraction of samples inside the two-sided $95\,\%$ interval of $\chi^{2}_{3}$, separately on the constant-velocity segments and inside the lane-change window.
\begin{enumerate}
\item[(a)] For which values is the filter overconfident, and for which is it underconfident? Use the NIS, not the RMSE, to answer.
\item[(b)] Explain why the RMSE on the constant-velocity segments first decreases and then increases with $\sigma_a$, whereas the RMSE inside the lane change decreases monotonically.
\item[(c)] Is there a value of $\sigma_a$ for which the NIS is consistent in both segments? Relate your answer to the closing remark of Example~\ref{ex:ekf-radar-cutin} on interacting multiple-model filters.
\end{enumerate}
A self-contained script and notebook are provided in the companion repository, in \repopath{companion\_\allowbreak repo/\allowbreak exercises/\allowbreak ch03/\allowbreak solutions/}.
\end{prob}

% ============================================
